\documentclass{article}
\usepackage{decomp}
\begin{document}

\pdfpage{1}
\noindent{\footnotesize bioRxiv preprint doi: https://doi.org/10.1101/134882; this version posted May 7, 2017. The copyright holder for this preprint (which was not certified by peer review) is the author/funder, who has granted bioRxiv a license to display the preprint in perpetuity. It is made available under a CC-BY 4.0 International license.}
\decompnote{This bioRxiv stamp is printed at the top of every one of the 47 pages; it is given once here. The PDF is the bioRxiv preprint (version posted May 7, 2017), not the later journal version.}

\begin{center}
{\Large\bfseries Whole-brain serial-section electron microscopy in larval zebrafish}
\end{center}

\noindent\textbf{David Grant Colburn Hildebrand\textsuperscript{1,2,3,4,5,†,*}, Marcelo Cicconet\textsuperscript{5}, Russel Miguel Torres\textsuperscript{2,4,†}, Woohyuk Choi\textsuperscript{6}, Tran Minh Quan\textsuperscript{6}, Jungmin Moon\textsuperscript{6}, Arthur Willis Wetzel\textsuperscript{7}, Andrew Scott Champion\textsuperscript{8}, Brett Jesse Graham\textsuperscript{4}, Owen Randlett\textsuperscript{2,†}, George Scott Plummer\textsuperscript{2,†}, Ruben Portugues\textsuperscript{2,†}, Isaac Henry Bianco\textsuperscript{2,†}, Stephan Saalfeld\textsuperscript{8}, Alex Baden\textsuperscript{9}, Kunal Lillaney\textsuperscript{9}, Randal Burns\textsuperscript{9}, Joshua Tzvi Vogelstein\textsuperscript{10}, Alexander Franz Schier\textsuperscript{2,3,11,12,13}, Wei-Chung Allen Lee\textsuperscript{4}, Won-Ki Jeong\textsuperscript{6}, Jeff William Lichtman\textsuperscript{2,3,14}, Florian Engert\textsuperscript{2,3,14,*}}

\begin{hangpars}\small
\textsuperscript{1} Graduate Program in Neuroscience, Division of Medical Sciences, Graduate School of Arts and Sciences, Harvard University, Cambridge, Massachusetts, United States of America

\textsuperscript{2} Department of Molecular and Cell Biology, Harvard University, Cambridge, Massachusetts, United States of America

\textsuperscript{3} Center for Brain Science, Harvard University, Cambridge, Massachusetts, United States of America

\textsuperscript{4} Department of Neurobiology, Harvard Medical School, Boston, Massachusetts, United States of America

\textsuperscript{5} Image and Data Analysis Core, Harvard Medical School, Boston, Massachusetts, United States of America

\textsuperscript{6} School of Electrical and Computer Engineering, Ulsan National Institute of Science and Technology (UNIST), Republic of Korea

\textsuperscript{7} Pittsburgh Supercomputing Center, Carnegie Mellon University, Pittsburgh, Pennsylvania, United States of America

\textsuperscript{8} Janelia Research Campus, Howard Hughes Medical Institute, Ashburn, Virginia, United States of America

\textsuperscript{9} Department of Computer Science, Johns Hopkins University, Baltimore, Maryland, United States of America

\textsuperscript{10} Department of Biomedical Engineering and Institute for Computational Medicine, Johns Hopkins University, Baltimore, Maryland, United States of America

\textsuperscript{11} Broad Institute of MIT and Harvard, Cambridge, Massachusetts, United States of America

\textsuperscript{12} Harvard Stem Cell Institute, Harvard University, Cambridge, Massachusetts, United States of America

\textsuperscript{13} FAS Center for Systems Biology, Harvard University, Cambridge, Massachusetts, United States of America

\textsuperscript{14} These authors contributed equally to this work.
\end{hangpars}

\begin{hangpars}\small
\textsuperscript{*}\textbf{Corresponding authors:}

D.G.C.H.: e-mail: david@hildebrand.name; telephone: +1 (857) 600-1770

F.A.E.: e-mail: florian@mcb.harvard.edu; telephone: +1 (617) 384-9773

\textsuperscript{†}\textbf{Present addresses:}

D.G.C.H.: Laboratory of Neural Systems, Rockefeller University, New York, New York, United States of America

R.M.T.: Allen Institute for Brain Science, Seattle, Washington, United States of America

G.S.P.: Tufts University School of Medicine, Boston, Massachusetts, United States of America

O.R.: Department of Cell and Developmental Biology, University of Pennsylvania Perelman School of Medicine, Philadelphia, Pennsylvania, United States of America

R.P.: Max Planck Institute of Neurobiology, Martinsried, Germany

I.H.B.: Department of Neuroscience, Physiology, and Pharmacology, University College London, London, United Kingdom
\end{hangpars}

\pdfpage{2}
Investigating the dense meshwork of wires and synapses that form neuronal circuits is possible with the high resolution of serial-section electron microscopy (ssEM)\textsuperscript{1}. However, the imaging scale required to comprehensively reconstruct axons and dendrites is more than 10 orders of magnitude smaller than the spatial extents occupied by networks of interconnected neurons\textsuperscript{2}—some of which span nearly the entire brain. The difficulties in generating and handling data for relatively large volumes at nanoscale resolution has thus restricted all studies in vertebrates to neuron fragments, thereby hindering investigations of complete circuits. These efforts were transformed by recent advances in computing, sample handling, and imaging techniques\textsuperscript{1}, but examining entire brains at high resolution remains a challenge. Here we present ssEM data for a complete 5.5 days post-fertilisation larval zebrafish brain. Our approach utilizes multiple rounds of targeted imaging at different scales to reduce acquisition time and data management. The resulting dataset can be analysed to reconstruct neuronal processes, allowing us to, for example, survey all the myelinated axons (the projectome). Further, our reconstructions enabled us to investigate the precise projections of neurons and their contralateral counterparts. In particular, we observed that myelinated axons of reticulospinal and lateral line afferent neurons exhibit remarkable bilateral symmetry. Additionally, we found that fasciculated reticulospinal axons maintain the same neighbour relations throughout the extent of their projections. Furthermore, we use the dataset to set the stage for whole-brain comparisons of structure and function by co-registering functional reference atlases and \textit{in vivo} two-photon fluorescence microscopy data from the same specimen. We provide the complete dataset and reconstructions as an open-access resource for neurobiologists and others interested in the ultrastructure of the larval zebrafish.

Pioneering studies in invertebrates established that synaptic-resolution wiring diagrams of complete neuronal circuits are valuable tools for relating a nervous system’s structure to its function\textsuperscript{3-6}. Such resources can be combined with perturbations, activity maps, or behavioural assays to examine how signalling through neuronal networks transforms information from the environment into relevant motor outputs\textsuperscript{5-10}. These studies benefited from the small size of the model organisms and stereotypy across individuals, which allow for complete ssEM of an entire individual or mosaicking of data from multiple individuals.

Vertebrate model nervous systems, on the other hand, are considerably larger and more variable. Consequently, ssEM of whole vertebrate neuronal circuits requires rapid computer-based technologies for acquiring, storing, and analysing many images from one animal. In many cases, anatomical data must be combined with other experiments on the same individual\textsuperscript{11-13} to define the relationship between structure, function, and behaviour. Because vertebrate nervous systems can vary substantially between individuals\textsuperscript{14}, distilling a particular circuit’s structure may require its full description in multiple individuals. Applied to mammalian brains, this kind of analysis would require imaging very large volumes that are technically still out of reach (but see ref. 15). Thus, in mammals, ssEM has been confined to smaller and often partial reconstructions of circuits\textsuperscript{16-23}. One strategy for capturing the brain-wide circuit underpinnings of behaviours is to generate whole-brain datasets with nanoscale resolution in smaller animals.

The larval zebrafish is an ideal system for this endeavour. It is a near-transparent vertebrate model organism that offers convenient optical access to its entire nervous system. Indeed, recent work has taken advantage of whole-brain calcium imaging in larval zebrafish\textsuperscript{24-27} and also ssEM examination of specific subregions\textsuperscript{28-30}. Integrated with an established genetic toolkit and a variety of quantitative behavioural assays\textsuperscript{28}, it is an excellent model organism for
\pdfpage{3}investigating the neuronal basis of behaviour\textsuperscript{31}. In addition, its small size makes the larval zebrafish ideally suited for ssEM examination of even the complete brain.

Our goal was to develop a framework for ssEM of complete larval zebrafish brains at 5–7 days post-fertilisation (dpf), which coincides with the emergence of complex behaviours such as prey capture\textsuperscript{32} and predator avoidance\textsuperscript{33}. Attempts to apply existing methods to the larval zebrafish were impeded by two obstacles: First, skin and membranes covering the brain\textsuperscript{34} prevent high-quality tissue preservation by whole-fish immersion alone. Second, consistent ultrathin sectioning is difficult to achieve in samples containing heterogeneous tissues but imperative for reconstructing three-dimensional structure from a series of two-dimensional sections.

To preserve ultrastructure across the complete brain with minimal damage, we developed a fine dissection technique to remove skin and membranes from the dorsum (see Methods; Extended Data Fig. 1a–c). Following dissection, traditional EM protocols (Extended Data Fig. 1d–h) resulted in high-quality tissue preservation throughout the brain (Extended Data Fig. 1i).

Consistent sectioning perpendicular to the majority of axon and dendrite paths is preferable for ease and reliability in reconstructing neuronal morphology. Therefore, we sectioned perpendicular to the long (anterior-posterior) axis of one 5.5dpf larval zebrafish, which required ∼2.5× more sections than the horizontal orientation (dorsal-ventral). This was made possible by customising an automated tape-collecting ultramicrotome\textsuperscript{35} to expand its tape-carrying capacity (Extended Data Fig. 2a–c) such that all sections could be collected on one continuous tape stretch. Furthermore, sectioning in test samples revealed that errors and loss occurred primarily when the composition of the sample changed dramatically (e.g., the border of embedded tissue and empty resin). We overcame this sectioning difficulty by embedding the sample in a low-viscosity epoxy resin within a surrounding support tissue of mouse cerebral cortex (Extended Data Fig. 1g–h, Extended Data Fig. 2f). This approach generated a stable library of sections on silicon wafers (Extended Data Fig. 2d–e) that could be imaged multiple times and at different resolutions.

Overview images were acquired to survey all sections (Extended Data Fig. 3; Supplementary Video 1), resulting in a stack of 17963 images spanning 1.02×10\textsuperscript{10}µm\textsuperscript{3}, consisting of 3.01×10\textsuperscript{11} voxels, and occupying 310 gigabytes of computer storage. In total, 17963 × ∼60nm-thick sections were collected from 18207 attempted, with 244 lost (1.34\%; Extended Data Fig. 3d top) and 283 containing partial regions of the tissue (1.55\%; Extended Data Fig. 3d middle; Extended Data Fig. 5). There were no adjacent losses and 5 instances of adjacent lost-partial or partial-partial events (0.03\%; Extended Data Fig. 3d bottom). From visualizations of this low-resolution data, we confirmed that our approach enabled stable serial sectioning through a millimetre-long region spanning from myotome 7 to the anterior-most larval zebrafish structures—encompassing the anterior spinal cord and the entire brain (Extended Data Fig. 3b–c; Supplementary Video 2).

We next selected sub-regions within this imaging volume to capture areas of interest at higher resolutions using multi-scale imaging\textsuperscript{35,36} (Fig. 1). We first performed nearly isotropic EM imaging by setting lateral resolution to match section thickness over the anterior-most 16000 sections (Fig. 1a–f; Supplementary Video 3). All cells are labelled in ssEM, so this volume offers a dense picture of the fine anatomy across the anterior quarter of the larval zebrafish including brain, sensory organs (e.g., eyes, ears, and olfactory pits), and other tissues. Furthermore, this resolution of 56.4×56.4×60nm\textsuperscript{3}/vx is ∼500× greater than that afforded by diffraction-limited light microscopy. The imaged volume spanned 2.28×10\textsuperscript{8}µm\textsuperscript{3}, consisted of 1.12×10\textsuperscript{12} voxels, and
\pdfpage{4}occupied 2.4 terabytes (TB). In this data, one can reliably identify cell nuclei and track large-calibre myelinated axons (Fig. 1b–c,e–f; Supplementary Video 4). To further resolve densely packed neuronal structures, a third round of imaging at 18.8×18.8×60nm\textsuperscript{3}/vx was performed to generate a high-resolution atlas specifically of the brain. The resulting image volume (Fig. 1d,g–h) spanned 12546 sections, contained a volume of 5.49×10\textsuperscript{7}µm\textsuperscript{3}, consisted of 2.36×10\textsuperscript{12} voxels, and occupied 4.9TB. Additional acquisition at higher magnifications was used to further inspect regions of interest, to resolve finer axons and dendrites, and to identify synaptic connections between neurons (Fig 1i–k). All images across sections and scales were co-registered form a coherent multi-resolution ssEM dataset (Supplementary Video 3).

With a framework in place for producing whole-brain ssEM data (Extended Data Fig. 6), we next tested the feasibility of identifying individual neurons and fluorescently defined regions across imaging modalities, as done for smaller volumes previously\textsuperscript{11-13}. To accomplish this, we matched the positions of nuclei in the ssEM dataset with those in two-photon fluorescence of a genetically encoded calcium indicator acquired from the same animal. The common structural features between the two image sets enabled identification of the same neurons in both datasets (Extended Data Fig.7; Supplementary Video 5). Furthermore, this revealed the imaging conditions, labelling density, and structural tissue features necessary for reliable matching across imaging modalities. It was difficult in regions where fluorescence signal was low (Extended Data Fig. 7e), where many cells were packed closely together (Extended Data Fig. 7f), and where progenitors add new neurons in between light microscopy and preparation for ssEM (Extended Data Fig. 7g). Improving the light-level data with specific labelling of all nuclei and faster light-sheet or other imaging approaches should greatly improve the ease and accuracy of matching neuron identity across modalities. This ability to assign neuron identity across imaging modalities demonstrates proof-of-principle for the potential integration of rich neuronal activity maps with subsequent whole-brain structural examination of functionally characterized neurons and their networks. By following a similar process, we were also able to transform the Z-Brain reference atlas\textsuperscript{37} and Zebrafish Brain Browser\textsuperscript{38}, expandable open-source zebrafish references containing a variety of molecular labels, to seamlessly integrate into the multi-resolution ssEM dataset (Extended Data Fig. 6; Extended Data Fig. 8).

We next tested the general applicability of this dataset for reconstructions of selected neurons. First, we reconstructed a peripheral lateral line afferent neuron that innervates hair cells in a dorsal neuromast sensory organ and projects to the hindbrain with its soma residing in the posterior lateral line ganglion (Fig. 2a–e; Supplementary Video 6). By re-imaging this neuromast at higher resolution (4×4×60nm\textsuperscript{3}/vx), we identified synapses connecting hair cells and this primary sensory neuron (Fig. 2b–c) while improving reconstructions in unmyelinated regions in the periphery (Fig. 2d). We also reconstructed myelinated motor neurons, including several in the spinal cord (Fig. 2f), that directly contacted muscle. Further, myelinated neurons could also be identified and annotated within the brain. These reconstructions highlight the utility of multi-resolution ssEM for reassembling morphologies of various large-calibre myelinated neurons from sensory inputs, throughout the brain, and back to peripheral innervation of muscle.

This ability to reconstruct individual neurons throughout the entire brain allows for the generation of complete circuit descriptions, a necessity for identifying the pathways that underlie entire behavioural programs. To exploit this potential opportunity, we produced a ‘projectome’ consisting of all the myelinated axons within this larval zebrafish’s brain (Fig. 3a–b; Supplementary Video 7; see Methods). We reconstructed 2589 myelinated axon segments along
\pdfpage{5}with many of their somata and dendrites to yield 39.9cm of combined axon and dendrite length. We could easily follow 834 myelinated axons comprising 30.6cm back to their somata. For the remaining axon segments (9.3cm), it was more difficult to identify their somata because of long unmyelinated stretches. The longest individual axon reconstructed, a trigeminal ganglion sensory afferent, was followed for 1.2mm and extended from its sensory terminals in the anterior skin surface to its termination field in the brain.

The resulting projectome included 94 lateral line afferents that innervated 41 neuromasts through their ganglia to the hindbrain (Fig. 3c,e). Hindbrain arborisations of lateral line afferents were, as expected, topographically arranged such that those innervating more anterior neuromasts were ventrolateral to those innervating more posteriorly located neuromasts\textsuperscript{39} (Fig. 3e). These reconstructions reveal a striking bilateral symmetry in the lateral line system (Supplementary Video 8). Only one neuromast and its afferents did not have counterparts on the other side (Supplementary Video 8, orange sphere).

The projectome also included a significant fraction of the reticulospinal neuron population, which sends myelinated axons to the spinal cord from the midbrain and hindbrain. Similar to what we observed in the lateral line system, these neurons exhibit bilateral symmetry (Fig. 4a, red label). Additionally, we were able to identify individual reticulospinal neurons by their known positions and morphologies\textsuperscript{40,41} (Fig. 3d,f; Fig. 4a; Extended Data Fig. 9i; Supplementary Video 9). This afforded us an opportunity to precisely examine the extent of symmetry in this system.

From the projectome, we selected 44 identified reticulospinal neurons (22 on each side) whose myelinated axons form the medial longitudinal fasciculus (MLF)\textsuperscript{41} to quantitatively analyse the degree of bilateral symmetry in specific projections. We started by investigating whether the myelinated axons in the MLF of one hemisphere are symmetric in their shape and position to that of the homologous neurons on the contralateral side. Notably, computing a globally optimal pairwise assignment (Fig. 4b, asterisks) resulted in excellent matching of left-right homologs (Fig. 4b, diagonal) in all but one pair (Fig. 4b, red asterisks; Extended Data Fig. 9d).

In addition, we also inquired about the possibility of bilateral symmetry in the neighbour relations between myelinated axons within the nerve bundle. In cross-sections, we observed that axons appeared to be situated in relatively similar positions within the left and right bundle, suggesting that bilateral symmetry could also extend to the neighbour relations within the MLF (Fig. 4d, compare relative positions across line). We developed quantitative metrics for comparing the spatial relations between an identified axon and its neighbours within the bundle on one side of the brain versus the other side (see Methods; Fig. 4e,g−h; Supplementary Video 10). We selected 6 left-right pairs of identified neurons (12 axons in total; Fig. 4f) and calculated the positional relationships of all the 15 (i.e., $\sum_{1}^{5} n$) pairwise combinations between neighbouring axons in one bundle with respect to the reflected pair from the other side (Extended Data Fig. 9e−j). We found that the positional arrangements of these axons are mirror symmetric with their counterparts on the contralateral side. This symmetry in neighbour relations was evident over long stretches of the nerve bundles (Fig. 4f−j; Supplementary Video 10). Moreover, the axon neighbour relations returned to a state of symmetry even in the face of local perturbations such as the entry of a new axon into the bundle (Fig. 4f−i; Supplementary Video 10). Similar relationships were seen in the larger set of 44 reticulospinal axons (Extended Data Fig. 9i−k).


\pdfpage{6}
Although these axons originate from stereotyped positions in the midbrain and hindbrain, we expected that their positions would become progressively more scrambled in the same way neighbour relations of axons in peripheral mammalian nerves change with distance and show no mirror symmetry\textsuperscript{14}. Were that the case, the configuration of the axons in the two bundles should become considerably less symmetric as they travel further posterior toward the spinal cord. The fact that this kind of randomization in position does not occur raises the possibility that axon bundles preserve some positional information along their course. Such an idea has a precedent in the ribbon-shaped optic nerve in certain kinds of fish, but was thought to be a special feature of this structure\textsuperscript{42}.

This work suggests that stereotyped neighbour relations may be a general feature of fasciculated central nerve tracts. Interestingly, not all axons seemed to obey this symmetry as consistently as others (Fig. 4j). It is intriguing that in mammals—where axon branches reorganize in development via activity dependent mechanisms—such symmetry is not known to occur. Perhaps fish have both plastic and hardwired axons in their central nervous system whose development is guided by distinct mechanisms.

Here we demonstrate the feasibility of reconstructing neurons in whole-brain ssEM datasets acquired from larval zebrafish. We illustrate the utility of multi-scale imaging for reducing both imaging time and data storage requirements. We show that imaging of section libraries on wafers allows repeated re-imaging at higher resolutions, an option not available with block-face ssEM methods. Furthermore, continued development of faster EM technologies\textsuperscript{43} will hasten the re-imaging process and permit whole-brain studies in a fraction of the time required here.

Finally, the datasets presented here are not limited to analyses of the nervous system and can easily be extended to examine the structure of other organ systems (Extended Data Fig. 10), such as musculoskeletal, cardiac, intestinal, and pancreatic tissues. The data we generated can, thus, be used as an anterior larval zebrafish reference atlas and is available as an online resource (http://hildebrand16.neurodata.io/) for study by the scientific community.


\pdfpage{7}
\section*{References}
\begin{hangpars}

1. Briggman, K. L. \& Bock, D. D. Volume electron microscopy for neuronal circuit reconstruction. \textit{Curr. Opin. Neurobiol.} \textbf{22}, 154-161 (2012).

2. Lichtman, J. W. \& Denk, W. The big and the small: challenges of imaging the brain's circuits. \textit{Science} \textbf{334}, 618-623 (2011).

3. White, J. G., Southgate, E., Thomson, J. N. \& Brenner, S. The structure of the nervous system of the nematode \textit{Caenorhabditis elegans}. \textit{Phil. Trans. R. Soc. Lond. B} \textbf{314}, 1-340 (1986).

4. Jarrell, T. A.\textit{ et al.} The connectome of a decision-making neural network. \textit{Science} \textbf{337}, 437-444 (2012).

5. Randel, N.\textit{ et al.} Neuronal connectome of a sensory-motor circuit for visual navigation. \textit{eLife} \textbf{3}, e02730 (2014).

6. Ohyama, T.\textit{ et al.} A multilevel multimodal circuit enhances action selection in Drosophila. \textit{Nature} \textbf{520}, 633-639 (2015).

7. Chalfie, M.\textit{ et al.} The neural circuit for touch sensitivity in \textit{Caenorhabditis elegans}. \textit{J. Neurosci.} \textbf{5}, 956-964 (1985).

8. Gray, J. M., Hill, J. J. \& Bargmann, C. I. A circuit for navigation in \textit{Caenorhabditis elegans}. \textit{Proc. Natl. Acad. Sci. U.S.A.} \textbf{102}, 3184-3191 (2005).

9. Chalasani, S. H.\textit{ et al.} Dissecting a circuit for olfactory behaviour in \textit{Caenorhabditis elegans}. \textit{Nature} \textbf{450}, 63-70 (2007).

10. Macosko, E. Z.\textit{ et al.} A hub-and-spoke circuit drives pheromone attraction and social behaviour in \textit{C. elegans}. \textit{Nature} \textbf{458}, 1171-1175 (2009).

11. Bock, D. D.\textit{ et al.} Network anatomy and \textit{in vivo} physiology of visual cortical neurons. \textit{Nature} \textbf{471}, 177-182 (2011).

12. Briggman, K. L., Helmstaedter, M. \& Denk, W. Wiring specificity in the direction-selectivity circuit of the retina. \textit{Nature} \textbf{471}, 183-188 (2011).

13. Lee, W.-C. A.\textit{ et al.} Anatomy and function of an excitatory network in the visual cortex. \textit{Nature} (2016).

14. Lu, J., Tapia, J. C., White, O. L. \& Lichtman, J. W. The interscutularis muscle connectome. \textit{PLoS Biol.} \textbf{7} (2009).

15. Mikula, S. \& Denk, W. High-resolution whole-brain staining for electron microscopic circuit reconstruction. \textit{Nat. Methods} \textbf{12}, 541-546 (2015).

16. Sjöstrand, F. S. Ultrastructure of retinal rod synapses of the guinea pig eye as revealed by three-dimensional reconstructions from serial sections. \textit{J. Ultrastruct. Res.} \textbf{2}, 122-170 (1958).

17. Sterling, P. Microcircuitry of the cat retina. \textit{Annu. Rev. Neurosci.} \textbf{6}, 149-185 (1983).

18. Hamos, J. E., Van Horn, S. C., Raczkowski, D. \& Sherman, S. M. Synaptic circuits involving an individual retinogeniculate axon in the cat. \textit{J. Comp. Neurol.} \textbf{259}, 165-192 (1987).

19. Somogyi, P., Tamás, G., Lujan, R. \& Buhl, E. H. Salient features of synaptic organisation in the cerebral cortex. \textit{Brain Res. Brain Res. Rev.} \textbf{26}, 113-135 (1998).

20. Shepherd, G. M. G. \& Harris, K. M. Three-dimensional structure and composition of CA3→CA1 axons in rat hippocampal slices: implications for presynaptic connectivity and compartmentalization. \textit{J. Neurosci.} \textbf{18}, 8300-8310 (1998).

21. Mishchenko, Y.\textit{ et al.} Ultrastructural analysis of hippocampal neuropil from the connectomics perspective. \textit{Neuron} \textbf{67}, 1009-1020 (2010).


\pdfpage{8}
22. da Costa, N. M., Fürsinger, D. \& Martin, K. A. C. The synaptic organization of the claustral projection to the cat's visual cortex. \textit{J. Neurosci.} \textbf{30}, 13166-13170 (2010).

23. Li, J., Erisir, A. \& Cline, H. \textit{In vivo} time-lapse imaging and serial section electron microscopy reveal developmental synaptic rearrangements. \textit{Neuron} \textbf{69}, 273-286 (2011).

24. Ahrens, M. B., Orger, M. B., Robson, D. N., Li, J. M. \& Keller, P. J. Whole-brain functional imaging at cellular resolution using light-sheet microscopy. \textit{Nat. Methods} \textbf{10}, 413-420 (2013).

25. Panier, T.\textit{ et al.} Fast functional imaging of multiple brain regions in intact zebrafish larvae using selective plane illumination microscopy. \textit{Front. Neural Circuits} \textbf{7}, 65 (2013).

26. Portugues, R., Feierstein, C. E., Engert, F. \& Orger, M. B. Whole-brain activity maps reveal stereotyped, distributed networks for visuomotor behavior. \textit{Neuron} \textbf{81}, 1328-1343 (2014).

27. Prevedel, R.\textit{ et al.} Simultaneous whole-animal 3D imaging of neuronal activity using light-field microscopy. \textit{Nat. Methods} \textbf{11}, 727-730 (2014).

28. Friedrich, R. W., Jacobson, G. A. \& Zhu, P. Circuit neuroscience in zebrafish. \textit{Curr. Biol.} \textbf{20}, R371-R381 (2010).

29. Friedrich, R. W., Genoud, C. \& Wanner, A. A. Analyzing the structure and function of neuronal circuits in zebrafish. \textit{Front. Neural Circuits} \textbf{7}, 71 (2013).

30. Wanner, A. A., Genoud, C., Masudi, T., Siksou, L. \& Friedrich, R. W. Dense EM-based reconstruction of the interglomerular projectome in the zebrafish olfactory bulb. \textit{Nat. Neurosci.} \textbf{19} (2016).

31. Fetcho, J. R., Higashijima, S.-i. \& McLean, D. L. Zebrafish and motor control over the last decade. \textit{Brain Res. Rev.} \textbf{57}, 86-93 (2008).

32. Bianco, I. H., Kampff, A. R. \& Engert, F. Prey capture behavior evoked by simple visual stimuli in larval zebrafish. \textit{Front. Syst. Neurosci.} \textbf{5}, 101 (2011).

33. Dunn, T. W.\textit{ et al.} Neural circuits underlying visually evoked escapes in larval zebrafish. \textit{Neuron} \textbf{89}, doi:10.1016/j.neuron.2015.12.021 (2016).

34. Xiao, T. \& Baier, H. Lamina-specific axonal projections in the zebrafish tectum require the type IV collagen Dragnet. \textit{Nat. Neurosci.} \textbf{10}, 1529-1537 (2007).

35. Hayworth, K. J.\textit{ et al.} Imaging ATUM ultrathin section libraries with WaferMapper: a multi-scale approach to EM reconstruction of neural circuits. \textit{Front. Neural Circuits} \textbf{8}, 68 (2014).

36. Kasthuri, N.\textit{ et al.} Saturated reconstruction of a volume of neocortex. \textit{Cell} \textbf{162}, 648-661 (2015).

37. Randlett, O.\textit{ et al.} Whole-brain activity mapping onto a zebrafish brain atlas. \textit{Nat. Methods} \textbf{12}, 1039-1046 (2015).

38. Marquart, G. D.\textit{ et al.} A 3D searchable database of transgenic zebrafish Gal4 and Cre lines for functional neuroanatomy studies. \textit{Front. Neural Circuits} \textbf{9}, doi:10.3389/fncir.2015.00078 (2015).

39. Alexandre, D. \& Ghysen, A. Somatotopy of the lateral line projection in larval zebrafish. \textit{Proc. Natl. Acad. Sci. U.S.A.} \textbf{96}, 7558-7562 (1999).

40. Kimmel, C. B., Powell, S. L. \& Metcalfe, W. K. Brain neurons which project to the spinal cord in young larvae of the zebrafish. \textit{J. Comp. Neurol.} \textbf{205}, 112-127 (1982).

41. Metcalfe, W. K., Mendelson, B. \& Kimmel, C. B. Segmental homologies among reticulospinal neurons in the hindbrain of the zebrafish larva. \textit{J. Comp. Neurol.} \textbf{251}, doi:10.1002/cne.902510202 (1986).


\pdfpage{9}
42. Scholes, J. H. Nerve fibre topography in the retinal projection to the tectum. \textit{Nature} \textbf{278} (1979).

43. Eberle, A. L.\textit{ et al.} High-resolution, high-throughput imaging with a multibeam scanning electron microscope. \textit{Journal of Microscopy} \textbf{259}, 114-120 (2015).
\end{hangpars}


\pdfpage{10}
\section*{Methods}

\subsection*{Animal care}

Adult zebrafish (\textit{Danio rerio}) for breeding were maintained at 28ºC on a 14hr:10hr light:dark cycle following standard methods\textsuperscript{44}. The \textit{Tg(elavl3:GCaMP5G)a4598} fish line\textsuperscript{24} used in this study was of genotype \textit{elavl3}:GCaMP5G\textsuperscript{+/+}; nacre (\textit{mitfa}\textsuperscript{–/–}), conveying nearly pan-neuronal expression of the calcium indicator GCaMP5G\textsuperscript{45} and increased transparency due to the nacre mutation\textsuperscript{46}. The larval zebrafish described in this study were raised in filtered fish facility water\textsuperscript{44} until 5–7 days post-fertilization (dpf).

Mice from which support tissue was collected had been previously euthanized for unrelated experiments. Only unused, to-be-discarded tissue from these studies was harvested to serve as support tissue.

The Standing Committee on the Use of Animals in Research and Training of Harvard University approved all animal experiments.

\subsection*{Two-photon laser-scanning microscopy}

Larval zebrafish were immobilized by immersion in 1mg mL\textsuperscript{–1} α-bungarotoxin (Invitrogen) and mounted dorsum-up in 2\% low-melting-temperature agarose in a small dish containing a silicone base (Sylgard® 184, Dow Corning). Upon agarose hardening, E3 solution (5mM NaCl, 0.17mM KCl, 0.33mM CaCl\textsubscript{2}, and 0.33mM MgSO\textsubscript{4}) was added to the dish. \textit{In vivo } structural imaging of \textit{elavl3}-driven nearly pan-neuronal GCaMP5G signal was conducted with a custom-built two-photon microscope with a Ti:Sapphire laser (Mai Tai®, Spectra-Physics) excitation source tuned to 800nm. Frames with a 764.4×509.6µm\textsuperscript{2 }field-of-view size (1200×800px\textsuperscript{2}) were acquired at 1µm intervals (0.637×0.637×1µm\textsuperscript{3}vx\textsuperscript{–1}) at ∼1 Hz with a scan pattern of four evenly spaced, interlaced passes, similar to a previously reported scanning approach\textsuperscript{26}. A low-noise anatomical snapshot of complete brain fluorescence was captured in 300 planes, each the sum of 50 single frames. All light-based imaging was performed without any intentional stimulus presentation.

\subsection*{Dissection and EM tissue preparation}

Any aqueous solutions were prepared with water passed through a purification system (typically Arium 611VF, Sartorius Stedim Biotech).

Larval zebrafish embedded in low-melting-temperature agarose were introduced to a dissection solution (64mM NaCl, 2.9mM KCl, 10mM HEPES, 10mM glucose, 164mM sucrose, 1.2mM MgCl\textsubscript{2}, 2.1mM CaCl\textsubscript{2}, and pH 7.5; ref. \textsuperscript{47}) containing 0.02\% (w/v) tricaine mesylate (MS-222, Sigma-Aldrich). Flow of red blood cells through the vasculature was confirmed before beginning the dissection as an indicator of good health. A portion of agarose was removed to expose the dorsum from the posterior hindbrain to the anterior optic tectum. The dissection was initiated by puncturing the thin epithelial layer over the rhombencephalic ventricle above the hindbrain\textsuperscript{48} with a sharpened tungsten needle. Small incremental anterior-directed incisions were then made along the midline as close to the surface as possible until the brain was exposed from the hindbrain entry to the middle of the optic tectum (Extended Data Fig. 1). The majority of damage associated with this dissection was restricted to medial tectal proliferation zone progenitor cells\textsuperscript{49} that are unlikely to have integrated into functional neuronal circuits.

Dissections lasted 1–2min, upon which time the complete dish was immersed in 2.0\% formaldehyde and 2.5\% glutaraldehyde fixative solution (Electron Microscopy Sciences)
\pdfpage{11}overnight at room temperature. Following washes, larval zebrafish were cut out from the dish in a block of agarose with a scalpel and moved to a round-bottom microcentrifuge tube. Specimens were then incubated in post-fixation solution containing 1\% osmium tetroxide and 1.5\% potassium ferricyanide for 2hr, washed with water, washed with 0.05M maleate buffer (pH 5.15), and stained with 1\% uranyl acetate in maleate buffer overnight. During the subsequent wash step with maleate buffer, larval zebrafish were freed from the surrounding agarose block and moved to a new microcentrifuge tube. Next, specimens were washed with water, dehydrated with serial dilutions of acetonitrile in water (25\%, 50\%, 70\%, 70\%, 80\%, 90\%, 95\%, 100\%, 100\%, 100\%) for 10 minutes each, and infiltrated with serial dilutions of a diepoxyoctane-based low viscosity resin\textsuperscript{50} in acetonitrile (25\%, 50\%, 75\%, 100\%) for 1hr each. The samples were then embedded in the same resin and hardened for 2–3d at 60°C. Additional solution, washing, and timing details were described previously in a step-by-step protocol\textsuperscript{51}.

\subsection*{Serial sectioning}

Sections were continuously cut with a diamond knife affixed to an ultramicrotome (EM UC6, Leica) and collected with a customised automated tape-collecting ultramicrotome\textsuperscript{35} onto 8mm-wide and 50–75µm-thick tape (Kapton® polyimide film, DuPont). Sectioning restarts were occasionally required for three reasons: fine-tuning of tape positioning or settings is necessary at the beginning of a run; the ultramicrotome design is constrained by a cutting window depth range of ∼200µm; and diamond knives must be shifted after cutting several thousand sections to expose the sample to a fresh cutting edge before dulling impairs sectioning quality. When necessary, restarts were completed as quickly as possible (typically 1–2min) in order to minimize possible thermal, electrostatic, or other fluctuations. For the same reason, tape reels were fed continuously rather than being reloaded. This required modifying the automated tape-collecting ultramicrotome to extend the device’s main mounting plate and enlarge tape-carrying reels to accommodate the entire length needed (compare Extended Data Fig. 2a with Fig. 1E from ref. \textsuperscript{35}). This combination of fast restarts and continuous tape feeding was successful at maintaining a steady state across restarts.

We sectioned two larval zebrafish specimens and these represent the only two samples we have attempted to cut with the surrounding support tissue approach. The primary focus of this study was a 5.5dpf larval zebrafish sectioned with a 45º ultra diamond knife (Diatome) and a nominal sectioning thickness that averaged 60nm with a variable setting ranging from 50–70nm depending on sectioning consistency. Restarts for this sample occurred after sections 276, 3669, 6967, 10346, 12523, 12916, and 15956. Knife shifts occurred after sections 6967 and 12916. After sectioning, the tape was cut into segments with a razor blade between collected sections and adhered with double-sided conductive carbon adhesive tape (Ted Pella) to 4in-diameter silicon wafers (University Wafer), which served as an imaging substrate. In total, 17963 × ∼60nm-thick sections were spread across 80 wafers for this specimen. A thin layer of carbon was deposited onto each wafer to prevent charging during scanning EM imaging.

One potential limitation of the 5.5dpf larval zebrafish section series presented herein is the section thickness, as minimizing section thickness is an important factor in the success of axon and dendrite reconstructions\textsuperscript{1}. Small neuronal processes (on the order of the section thickness) are difficult to reconstruct in thicker sections especially when they are running roughly parallel to the plane of the section. To be sure that our approach was not fundamentally limited to thicker sections, we sectioned the second sample—a 7dpf larval zebrafish—with a nominal sectioning thickness that remained constant at 50nm throughout the entire cutting
\pdfpage{12}session using a 45º histo diamond knife (Diatome). Restarts for this sample occurred after sections 296, 312, 4114, 8233, and 12333. Knife shifts occurred after sections 4114 and 12333. In total, we obtained 15046 × ∼50nm-thick sections from 15052 attempted (Extended Data Fig. 5) and spread them across 70 wafers. The thinner sections did not cause more lost material: this series contained 6 losses (0.04\%; Extended Data Fig. 5d top), 25 partial sections (0.17\%; Extended Data Fig. 5d middle), no adjacent losses, and 6 adjacent lost-partial or partial-partial events (0.04\%; Extended Data Fig. 5d bottom).
\decompnote{The 7dpf series statistics (15046 sections, 6 lost, 25 partial, 6 adjacent events) are shown in Extended Data Fig. 4 (panel d); Extended Data Fig. 5 shows partial sections of the 5.5dpf series. "Extended Data Fig. 5" / "5d" here seem to mean Extended Data Fig. 4 / 4d.}

A nominal section thickness of ∼60nm made it possible to span the entire 5.5dpf larval zebrafish brain in ∼18000 sections, as determined by finding the location of the spinal cord-hindbrain boundary\textsuperscript{52}. The 7dpf sample was sectioned at 50nm, but was not made the focus of imaging because of an over-trimming error that caused less of the brain to be captured. However, improved reliability for this sample despite a ∼17\% reduction in nominal sectioning thickness suggests that higher axial resolution is readily attainable, though expanding the tape-carrying capacity of the microtome further may be necessary. A section thickness of ≤30nm would increase confidence in the ability to reconstruct complete neuronal circuit connectivity and this thickness will be the target in our future work. Thicknesses of ≤30nm are possible for mammalian brain sections of similar areas\textsuperscript{36,53}, suggesting that the approach described herein will not be a limiting factor.

Once wafers contained tape segments, they were made hydrophilic by glow discharging very briefly, post-section stained for 1–2min inside a chamber containing sodium hydroxide pellets with a stabilised lead citrate solution (Leica, UltroStain II) filtered through a 0.2µm syringe filter, and then washed thoroughly with boiled water.

\subsection*{Electron microscopy}

WaferMapper software was used with light-based wafer overview images to semi-automatically map the positions of all sections and relate them to per-wafer fiducial markers that enabled targeted section overview acquisition at a resolution of 758.8×758.8×60nm\textsuperscript{3}vx\textsuperscript{–1}. Semi-automated alignment of section overviews in WaferMapper then permitted targeting for higher resolution imaging\textsuperscript{35}.

Field emission scanning EM of back-scattered electrons was primarily conducted on a Zeiss Merlin equipped with a large-area imaging scan generator (Fibics) and stock back-scattered electron detector. An accelerating voltage of 5.0kV and beam current of 7–10nA was used for most acquisition. Imaging of back-scattered electrons at the highest resolutions (4×4×60nm\textsuperscript{3}vx\textsuperscript{–1}) was performed on an FEI Magellan XHR 400L with an accelerating voltage of 5.0kV and a beam current of 1.6–3.2A. Fields of view acquired on a given section varied depending on the area occupied by the sample in cross-section. All acquisition was performed with a scan rate at or under 1Mpx s\textsuperscript{–1}. This resulted in acquisition times of 5.4 days for section overview images (758.8×758.8×60nm\textsuperscript{3}vx\textsuperscript{–1}), 97 days for full transverse images (56.4×56.4×60nm\textsuperscript{3}vx\textsuperscript{–1}), and 100 days for high-resolution brain atlas images (18.8×18.8×60nm\textsuperscript{3}vx\textsuperscript{–1}).

\subsection*{Image alignment and intensity normalization}

Producing anatomically consistent registration over our relatively long section series required control of region-of-interest drift, over-fitting, magnification changes, and intensities. In order to quickly assess the quality of the dataset and begin reconstructions, we initially performed affine intra- and inter-section EM image registrations with Fiji\textsuperscript{54} TrakEM2 alignment plug-ins\textsuperscript{55}. These results revealed that additional nonlinear registration was required in order to
\pdfpage{13}compensate for distortions that were likely caused by section compression during cutting and minor sample charging during imaging. While the state-of-the-art elastic registration method\textsuperscript{56} also provided in Fiji\textsuperscript{54} as a TrakEM2 alignment plug-in was capable of achieving excellent local registration, we experienced difficulty—at least without modification to the existing implementation—in achieving an anatomically consistent result that preserved the overall structure of the larval zebrafish, largely due to our inability to successfully constrain region-of-interest drift and magnification changes. Furthermore, we determined that the similar AlignTK\textsuperscript{11} method, which uses Pearson correlation as the matching criterion coupled with spring mesh relaxation to stabilize the global volume, was likely to suffer from similar problems and would require substantial additional data handling to operate on our multi-resolution dataset.

In order to preserve the overall structure of the larval zebrafish and simultaneously achieve high-quality local registration, we turned to a new Signal Whitening Fourier Transform Image Registration (SWiFT-IR) method\textsuperscript{53,57}. Compared to conventional Pearson or phase correlation-based registration approaches, SWiFT-IR produces more robust image matchings by using modulated Fourier transform amplitudes, adjusting its spatial frequency response during the image matching steps to maximize a signal-to-noise measure that serves as its main indicator of alignment quality. This more reliable alignment signal in turn better handles variations in biological content and typical data distortions. Furthermore, SWiFT-IR achieves higher precision in block matching as a result of the signal whitening, improves computational speeds thanks to computational complexity advantages of the fast-Fourier transform, and reduces iterative convergence from thousands to dozens of steps. Together, these capabilities enable a model-driven alignment in place of the usual approach of comparing and aligning a given section to a pre-selected number of adjacent sections.

The SWiFT-IR model we used consisted of an estimate of local aligned volume content formed by a windowed average, typically spanning ∼6 µm along the z axis orthogonal to the sectioning plane. Any severely damaged regions, in particular partial sections, were removed from the model so that they would not adversely influence alignment results. This model then served as a template during registration, where raw images were matched to the current model rather than directly to a subset of nearby sections. Alignment proceeded in an iterative fashion starting at the lowest resolution acquired (section overviews, 758.8×758.8×60 nm\textsuperscript{3}vx\textsuperscript{–1}) and progressing incrementally to the highest resolution image set (56.4×56.4×60 nm\textsuperscript{3}vx\textsuperscript{–1} for most regions outside the brain, 18.8×18.8×60 nm\textsuperscript{3}vx\textsuperscript{–1} for most regions inside the brain).

At each resolution, the source images were iteratively aligned to the current model until no further significant gain in alignment could be achieved, as indicated by the SWiFT-IR signal-to-noise figure-of-merit. The model was then transferred to the next higher resolution data by applying the current warpings to source data for that scale. Iterative model refinement then continued at this subsequent resolution level. Although the majority of computations were locally affine, residual nonlinear deformations, particularly at the highest resolutions, were represented by a triangulation mesh that deformably mapped raw data onto the model volume.

Importantly, access to the lowest resolution section overview data for each section permitted us to build an initial model that constrained subsequent registration steps to the overall structure of the larval zebrafish. Although the resolution and signal quality of overview images were intentionally kept low in favour of rapid acquisition, the fact that overviews were all quickly captured with the same microscope settings and included surrounding support tissue provided key constraints for model refinement that resulted in a more accurate global result.


\pdfpage{14}
More specifically, the full 17963 section overview image set was processed using SWiFT-IR to produce an initial aligned overview model at 564×564×600 nm\textsuperscript{3}vx\textsuperscript{–1} resolution. Although the lowest resolution section overview images were each captured at 758.8×758.8×60 nm\textsuperscript{3}vx\textsuperscript{–1}, the relative oversampling in the stack direction enabled a geometrically accurate model at 564×564×600 nm\textsuperscript{3}vx\textsuperscript{–1} resolution. This initial model was then cropped to the bounds of the larval zebrafish and warped, using SWiFT-IR–driven matching across the midline axis, to remove cutting compression, rotations, and other systematic variations in the pose of the specimen. The 16000 section 56.4×56.4×60 nm\textsuperscript{3}vx\textsuperscript{–1} image set was next downsampled to 564×564×600 nm\textsuperscript{3}vx\textsuperscript{–1} to match the overview model. It was then aligned to the overview model, resulting in an improved model. The matching and remodeling process was iterated at this scale until there was no further improvement in the SWiFT-IR match quality. The final model at this scale was then expanded to 282×282×300 nm\textsuperscript{3}vx\textsuperscript{–1}, or one-sixth scale, and similarly aligned in an iterative fashion until there was no further improvement. This one-sixth scale model volume (∼6 Gvx, 1600×1400×2667 vx) was convenient for rapid viewing to identify and manually correct defects and refine the pose. Further scales at 169.2×169.2×180 nm\textsuperscript{3}vx\textsuperscript{–1} and 56.4×56.4×60 nm\textsuperscript{3}vx\textsuperscript{–1} were similarly processed by successively expanding the model and aligning at each resolution until there was no further significant improvement in the figure-of-merit. Finally, the 12546-section highest resolution 18.8×18.8×60nm\textsuperscript{3}vx\textsuperscript{–1} image set was similarly registered using the final 56.4×56.4×60nm\textsuperscript{3}vx\textsuperscript{–1} volume as its model.

Image intensity was adjusted across sections for each dataset to achieve a consistent background level chosen as the average over a tissue-free region, typically a 256×256px\textsuperscript{2} area above and to the right of the larval zebrafish. Many images were acquired at 16-bit depth and were converted in this process to 8-bit depth. The target background level was mapped to intensity 250, which left headroom for bright pixels while keeping most tissue of interest from saturating the range. Next, a linear intensity fit between the background and a second level, typically the average grey level of a variable size continuous trajectory region on the right side of the brain, was made to adjust the intensity values for each section.

More information on SWiFT-IR software tools is publicly available online (http://www.mmbios.org/swift-ir-home).

\subsection*{Image annotation and neuron reconstructions}

Reconstructing neuronal structures across multi-resolution ssEM image volumes acquired from the same specimen profits from being able to simultaneously access and view separate but co-registered datasets. Without this ability, some of the time benefits of our imaging approach would be offset by the need to register and track structures across volumes that span both low-resolution, large fields of view and high-resolution, specific regions of interest. With this in mind, we added a feature to the Collaborative Annotation Toolkit for Massive Amounts of Image Data (CATMAID) neuronal circuit mapping software\textsuperscript{58,59} to overlay and combine image stacks acquired with varying resolutions in a single viewer (Extended Data Fig. 6). Multiple image stacks are rendered using WebGL, which makes it possible to present co-registered stacks of different resolutions in the same view. Furthermore, this feature combines stacks via a configurable overlay order, introduces blending operations for each overlaid stack, and enables programmatic shaders for dynamic image processing. When resolutions of overlaid stacks differ, the nearest available zoom level for each stack is interpolated either bilinearly or via nearest neighbours. Missing regions in each stack can be either omitted or rendered with nearest neighbour interpolation. To account for the increase in data storage and bandwidth when viewing
\pdfpage{15}multiple image stacks, the CATMAID image data hierarchy was extended with a shared graphics card memory cache of image tiles using a least-recently-used replacement policy. All additions and modifications to the CATMAID software are now incorporated into the main release version (http://github.com/catmaid/CATMAID).

Manual reconstruction was conducted using our modified version of CATMAID by placing nodes near the centre of each neuronal structure on every section in which it could be clearly identified. This led to a wire-frame model of each annotated structure. Starting points for reconstruction (“seeds”) of myelinated processes were manually identified by searching all tissue on a given section from top-left to bottom-right at the highest available resolution for profiles surrounded by the characteristic thick, densely stained outline associated with staining of the myelin sheath\textsuperscript{60} (see Fig. 1e,g,i). To obtain seeds for reconstruction of the complete projectome for one larval zebrafish, this process was repeated every 50 sections throughout the set of 16000 imaged at 56.4×56.4×60nm\textsuperscript{3}vx\textsuperscript{–1} resolution. Many such reconstructions were produced in an initial, affine-only image alignment space before being mapped into the final, nonlinear image alignment space. The reconstructions reported here represent ∼450 days of human annotation.

For visualization and reported length measurements, each mapped wire-frame was smoothed using custom python-based software that implemented a Kalman smoothing algorithm on a space defined by manually annotated points within unique segments. Initial state variables for smoothing were derived by an optimization of point-to-line distance to connected reconstructed segments. Other variables were tuned with the Estimation Maximization algorithm of the pykalman library to compensate for a lack of human input where data was unavailable due to lost or partial sections. While the final image alignment was high quality, smoothing the reconstructions in this manner should produce a slight underestimate rather than an overestimate of reported reconstruction path lengths.

Neurons with known projection patterns or identities were named accordingly in the CATMAID database. For example, the reconstruction of a neuron innervating the right anterior macula (utricle) was named as “Ear\_AnteriorMacula\_R\_01”, while an identified neuron such as the left Mauthner neuron was named “Mauthner\_L”. Two identifiable neurons on each side belonged to the “MeM” class, which emanates from the nucleus of the medial longitudinal fasciculus. On each side, these were differentiated into dorsal (MeMd) and ventral (MeMv) subclasses based on consistent soma positioning.

\subsection*{Correspondence across light and EM datasets}

Correspondence of individual neurons or, in the case of the Z-Brain atlas, small regions across imaging modalities was achieved with landmark-based three-dimensional (3D) thin-plate spline warping of each fluorescence dataset to the ssEM dataset using BigWarp\textsuperscript{61}. For matching \textit{in vivo} two-photon laser-scanning microscopy data to ssEM data from the same specimen, we primarily chose landmarks consisting of distinctive arrangements of low-fluorescence regions where GCaMP5G was excluded and could be easily matched to similar patterns of neuronal nuclei at the EM level. For matching the Z-Brain atlas\textsuperscript{37} to the ssEM dataset, we chose landmarks based on identifiable structures—such the region boundaries, known clusters of neurons, midline points, ganglia, and the outline of the brain—in the Z-Brain averaged \textit{elavl3}:H2B-RFP or anti-tERK fluorescence image stacks that were also observed in the ssEM data. The same Z-Brain landmarks were used for transforming a version\textsuperscript{62} of the Zebrafish Brain Browser\textsuperscript{38} that was previously registered into the Z-Brain atlas.

\subsection*{Data and reconstruction visualizations}


\pdfpage{16}
The image volumes and reconstructions were visualized using Vivaldi\textsuperscript{63}, a domain-specific language for rendering and processing on distributed systems, because it provides access to the parallel computing power of multi-GPU systems with language syntax similar to python.

For volume visualizations, we used a direct volume rendering ray-casting technique in which an orthogonal or perspective iterator was marched along a viewing ray while sampled voxel colours were accumulated using an alpha compositing algorithm. We screened out regions containing only support tissue during rendering with labelled volumes constructed by interpolating between manually produced masks indicating which image voxels belong to each separate tissue region. In cases where separate image volumes of the same region were rendered together (such as when ssEM data and fluorescence data were combined), direct volume rendering was performed by combining front-to-back colour and alpha compositions formed from the different transfer function belonging to each image volume.

For volume visualizations that included reconstructions, direct volume rendering of image data was combined with streamline rendering of reconstructions using two different techniques. The first combined an OpenGL framebuffer with the Vivaldi volume rendering. In this case, each streamline was rendered using OpenGL as a tube into an off-screen buffer (i.e., Framebuffer Object). Vivaldi then compared the resulting render and depth buffers to perform direct volume rendering of only the image data above the streamline depth value. This made it possible to ignore any image data voxels located behind streamlines, which were treated as opaque. The second technique involved generating a complete streamline volume by 3D rasterization. This streamline volume was then combined with the image data volume for direct volume rendering. The former technique is faster and can cope with dynamic streamline changes, but the latter was found to yield better overall rendering quality. Visualizations of reconstructions without the image volume context were similarly rendered as streamlines in Vivaldi or in the CATMAID 3D WebGL viewer.

Reference plane (e.g., horizontal, sagittal, and section) images were rendered by detecting the zero-crossing of each viewing ray and the plane. Support for viewing these opaque data views in some spatial regions alongside the semi-transparent volume visualization views in other regions was introduced as a new function (clipping\_plane) in Vivaldi. Similarly, contour (nonplanar) reslice support was developed to illustrate a flattened view along a specific contour (or neuron path) consisting of vertical line segments extracted from image data.

For many cases, the size of the image volume being rendered was larger than the available memory. In order to support out-of-core processing, we developed and integrated into Vivaldi a slice-based streaming computing framework using Hadoop distributed file system (HDFS) that will be reported elsewhere.

\subsection*{Symmetry analyses}

Initial observations of apparent symmetry in subsets of myelinated axon reconstructions were found during visual inspection (Fig. 3; Supplementary Video 8; Supplementary Video 9). To more quantitatively assess the extent of symmetry, we developed methods for 3D symmetry plane fitting and two symmetry analyses: one that produces a cost associated with the shape and position similarity between reconstructed 3D structures and another that compares the relative two-dimensional (2D) positioning of identified reconstruction pairs on one side with the relative positioning for the same identified pair on the opposite side at a given cross-section level (Fig. 4d). Only the longest reconstructed path from the soma through the myelinated axon projection was considered in plane fitting and subsequent analyses. Reconstructed dendrites or short axonal
\pdfpage{17}branches were ignored. Each resulting reconstruction path (skeleton) was represented as an ordered list of nodes (points) taken directly from manual reconstructions. Sidedness (left/right) for each reconstruction was determined by its soma position.

Most of these new analyses approaches, including symmetry plane fitting and 3D symmetry comparisons, are described elsewhere\textsuperscript{64}. The symmetry plane fitting, in brief, involves choosing an approximate symmetry plane, reflecting the complete set of points belonging to the reconstruction subset of interest with respect to this plane, registering the original and reflected point clouds with an iterative closest point algorithm, and inferring the optimal symmetry plane from the reflection and registration mappings. The subset of reconstructions from which this plane fitting was performed consisted of all identified neurons whose axon projections formed part of the medial longitudinal fasciculus (MLF), determined primarily using refs. \textsuperscript{40,41}.

The 3D symmetry comparison for each given template reconstruction on one side, in brief, involves reflecting all reconstructions from the opposite side and computing a matching cost via dynamic time warping (DTW) for each comparison. For our purposes, the DTW cost is taken as the sum of the Euclidian distances between all matched points normalized to the number of matched point pairs (Extended Data Fig. 9a−c). The DTW gap cost parameter for matching a point in one sequence with a gap in another is set to zero because our data is sampled at a nearly constant rate and we seek the optimal subsequence match even in cases where one sequence is shorter than or offset with respect to the other. To compensate for unmatched regions (i.e., overhangs), the DTW cost is then multiplied by a penalty factor that is proportional to the lengths of the sequences that remain unmatched (total length divided by matched length). Comparing each reconstruction on one side to all reconstructions from the opposite side formed a cost matrix from which an optimal pairwise assignment could be determined without any bias introduced from the previously determined identities. The Munkres algorithm\textsuperscript{65} (also known as the Hungrarian method) was used to compute the globally optimal pairwise assignment between reconstructions. For visualization clarity and to decrease analysis time, the subset of reconstructions to which this 3D symmetry analysis was applied was restricted to identified neuron classes with only 1-2 members per side whose axon projections formed part of the MLF, determined primarily using ref. \textsuperscript{41}.

We next sought to compare the relative 2D positioning of axon pairs with the relative 2D positioning of the same pair on the opposite side at a given cross-section level (Fig. 4d). For visualization clarity, the subset of reconstructions to which this relative 2D positioning analysis was applied was restricted to the Mauthner cell and the identified neurons whose somata resided in the nucleus of the medial longitudinal fasciculus.

First, a small angle offset in the sectioning plane relative to the true transverse plane was compensated for by projecting the point coordinates of reconstructions such that the previously computed symmetry plane became the plane \textit{x} = 0. Given the transverse planes $z_0$, $z_1$, and a projected \textit{skeleton S} containing points \textit{s} = (\textit{s}\textsubscript{\textit{x}}, \textit{s}\textsubscript{\textit{y}}, \textit{s}\textsubscript{\textit{z}}), we let $S_{z_0,z_1} = \{s \in S : z_0 \leq s_z < z_1\}$. That is, $S_{z_0,z_1}$ was taken as the subset of points from \textit{S} whose coordinates $s_z$ are all in the interval $[z_0, z_1)$. We refer to the subset of ℝ\textsuperscript{3} bounded by $z \in [z_0, z_1)$ as the \textit{slice} $[z_0, z_1)$. For each slice $[z_0, z_1)$ and skeleton \textit{S}, we defined $\langle S_{z_0,z_1} \rangle$ as the mean of the elements in $S_{z_0,z_1}$. This mean was then taken as representative of the skeleton \textit{S} in slice $[z_0, z_1)$ both for further analysis and plotting. Note that all analysis and plotting presented in static form (Fig. 4) was based on a slice thickness corresponding to a single section (∼60nm), in which case each slice consisted simply of each pair of adjacent sections. This was not the case for dynamic presentation (Supplementary Video 10) due to video time and size constraints.


\pdfpage{18}
For comparing one pair with its counterpart (Extended Data Fig. 9e), we then took \textit{s}\textsubscript{1}, ..., \textit{s}\textsubscript{\textit{n}} to be the set of representative points in a fixed slice for skeletons \textit{S}\textsubscript{1}, ..., \textit{S}\textsubscript{\textit{n}} and took \textit{t}\textsubscript{1}, ..., \textit{t}\textsubscript{\textit{n}} to be the representative points (for the same slice) of the respective skeletons \textit{T}\textsubscript{1}, ..., \textit{T}\textsubscript{\textit{n}} that were previously paired to \textit{S}\textsubscript{1}, ..., \textit{S}\textsubscript{\textit{n}} by the Munkres algorithm assignment in the 3D symmetry analysis, which for this subset of reconstructions always implied assignment of each identified neuron to its homolog on the opposite side.

To quantify the degree of similarity between any given pair and its opposite side counterpart, we devised two measures (Extended Data Fig. 9e). The first, termed the \textit{angle difference}, \textit{a}\textsubscript{\textit{i,j}}, between pairs, is defined as:
\[ a_{i,j} = \frac{1}{2}\left(1 - \frac{\langle s_j - s_i, t_j^s - t_i^s \rangle}{\lVert s_j - s_i \rVert \, \lVert t_j^s - t_i^s \rVert}\right). \]
The second, termed the \textit{distance difference}, \textit{d}\textsubscript{\textit{i,j}}, between pairs, is defined as:
\[ d_{i,j} = \frac{\bigl|\, \lVert s_j - s_i \rVert - \lVert t_j^s - t_i^s \rVert \,\bigr|}{M}, \]
where \textit{i}, \textit{j} are skeleton indices, $\{t_i^s\}$ is the reflections of $\{t_i\}$ with respect to the computed plane of symmetry, and \textit{M} is the maximum of $\bigl|\, \lVert s_j - s_i \rVert - \lVert t_j^s - t_i^s \rVert \,\bigr|$ across all pairs and all stacks. Note that \textit{a}\textsubscript{\textit{i,j}} and \textit{d}\textsubscript{\textit{i,j}} are normalized such that they can vary from 0 (no difference, 0° or 0µm) to 1 (maximum difference, 180° or ∼8µm). Further, when the points \textit{s}\textsubscript{\textit{i}} and \textit{s}\textsubscript{\textit{j}} are perfectly symmetric with respect to points \textit{t}\textsubscript{\textit{i}} and \textit{t}\textsubscript{\textit{j}}, then \textit{a}\textsubscript{\textit{i,j}} = 0 and \textit{d}\textsubscript{\textit{i,j}} = 0.

To more easily visualize this quantification, a \textit{difference matrix}, \textit{D}, was generated for each slice such that \textit{D}(\textit{i},\textit{j}) = \textit{a}\textsubscript{\textit{i,j}} if \textit{j} > \textit{i} and \textit{D}(\textit{i},\textit{j}) = \textit{d}\textsubscript{\textit{i,j}} if \textit{j} < \textit{i} (Fig. 4e; Extended Data Fig. 9f; Supplementary Video 10). Calculating the variance for each element in \textit{D} across all slices showed which pairs deviated most with respect to the opposite side counterpart (Fig. 4j). Heatmaps of vectorised (linearized) upper (\textit{j} > \textit{i}) and lower (\textit{j} < \textit{i}) triangles of \textit{D} across slices with the same pair ordering additionally revealed where differences were present for single pairs and their opposite side counterparts (Fig. 4g−h; Extended Data Fig. 9g−h; Supplementary Video 10). Plotting the sum of all \textit{a}\textsubscript{\textit{i,j}} and \textit{d}\textsubscript{\textit{i,j}} values for a given slice further illustrated where differences were present across the entire subset (Fig. 4i). Finally, the exact same process was performed after swapping the identity (assignment) of the two reconstructions with the lowest 3D symmetry analysis costs (MeLc and Mauthner neurons) to provide a basis for comparison (Fig. 4i).

\subsection*{Online data hosting}

All aligned ssEM data, reconstructions, co-registered Z-Brain data, and a guide outlining ways to interact with the data are publicly available online (http://hildebrand16.neurodata.io/). Our enhanced version of CATMAID\textsuperscript{58,59} software with multi-resolution support is used for visualization (Extended Data Fig. 6). All image data is served as a series of 8-bit 1024×1024px\textsuperscript{2} PNG images with a tRNS value of 255 specified to enable transparency. To provide a smooth visualization experience, the original resolution for each image stack was down-sampled multiple times to create a resolution hierarchy, where each level in the hierarchy corresponds to an image that is half the size as in the previous level. The entire down-sampled dataset required 2.8 TB of disk space (608 GB for 56.4×56.4×60nm\textsuperscript{3}vx\textsuperscript{–1} resolution EM data, 1.8 TB for 18.8×18.8×60nm\textsuperscript{3}vx\textsuperscript{–1} resolution EM data, 356 GB for 4×4×60nm\textsuperscript{3}vx\textsuperscript{–1} resolution of dorsal neuromasts, and 29 GB for 300×300×300nm\textsuperscript{3}vx\textsuperscript{–1} resolution of Z-Brain data). To serve data to the end user, we used the Amazon Web Services (AWS) Simple Storage Service (S3). The AWS
\pdfpage{19}S3 built-in web server was used to serve the static image files. A public CATMAID instance was then deployed on Amazon EC2 to point to the images in S3.

Custom tools generated for data handling and analysis are publicly available (http://github.com/davidhildebrand/hildebrand16/).


\pdfpage{20}
\section*{Methods References}
\begin{hangpars}

44. Brand, M., Granato, M. \& Nüsslein-Volhard, C. in \textit{Zebrafish: A Practical Approach} (eds Nüsslein-Volhard, C. \& Dahm, R.) Ch. 1, 7-37 (Oxford, 2002).

45. Akerboom, J.\textit{ et al.} Optimization of a GCaMP calcium indicator for neural activity imaging. \textit{J. Neurosci.} \textbf{32}, 13819-13840 (2012).

46. Lister, J. A., Robertson, C. P., Lepage, T., Johnson, S. L. \& Raible, D. W. nacre encodes a zebrafish microphthalmia-related protein that regulates neural-crest-derived pigment cell fate. \textit{Development} \textbf{126} (1999).

47. Ma, L.-H., Gilland, E., Bass, A. H. \& Baker, R. Ancestry of motor innervation to pectoral fin and forelimb. \textit{Nature Communications} \textbf{1}, 49 (2010).

48. Turner, M. H., Ullmann, J. F. P. \& Kay, A. R. A method for detecting molecular transport within the cerebral ventricles of live zebrafish (Danio rerio) larvae. \textit{J. Physiol.} \textbf{590}, 2233-2240 (2012).

49. Mueller, T. \& Wulliman, M. F. \textit{Atlas of early zebrafish brain development: a tool for molecular neurogenetics}. 1 edn, 183 (Elsevier, 2005).

50. Luft, J. H. in \textit{Advanced Techniques in Biological Electron Microscopy} (ed Koehler, J. K.) Ch. 1, 1-34 (Springer-Verlag, 1973).

51. Hildebrand, D. G. C. \textit{Whole-brain functional and structural examination in larval zebrafish}, Harvard University, (2015).

52. Ma, L.-H., Punnamoottil, B., Rinkwitz, S. \& Baker, R. Mosaic hoxb4a neuronal pleiotropism in zebrafish caudal hindbrain. \textit{PLoS ONE} \textbf{4}, e5944 (2009).

53. Morgan, J. L., Berger, D. R., Wetzel, A. W. \& Lichtman, J. W. The fuzzy logic of network connectivity in mouse visual thalamus. \textit{Cell} \textbf{165}, 192-206 (2016).

54. Schindelin, J.\textit{ et al.} Fiji: an open-source platform for biological-image analysis. \textit{Nat. Methods} \textbf{9}, 676-682 (2012).

55. Saalfeld, S., Cardona, A., Hartenstein, V. \& Tomančák, P. As-rigid-as-possible mosaicking and serial section registration of large ssTEM datasets. \textit{Bioinformatics} \textbf{26}, i57-i63 (2010).

56. Saalfeld, S., Fetter, R., Cardona, A. \& Tomancak, P. Elastic volume reconstruction from series of ultra-thin microscopy sections. \textit{Nat. Methods} \textbf{9} (2012).

57. Wetzel, A. W.\textit{ et al.} Registering large volume serial-section electron microscopy image sets for neural circuit reconstruction using FFT signal whitening. \textit{Applied Imagery Pattern Recognition Workshop (AIPR) 2016 IEEE} \textbf{in press} (2016).

58. Saalfeld, S., Cardona, A., Hartenstein, V. \& Tomančák, P. CATMAID: collaborative annotation toolkit for massive amounts of image data. \textit{Bioinformatics} \textbf{25}, 1984-1986 (2009).

59. Schneider-Mizell, C. M.\textit{ et al.} Quantitative neuroanatomy for connectomics in Drosophila. \textit{eLife} \textbf{5}, e12059 (2016).

60. Peters, A., Palay, S. L. \& Webster, H. d. The fine structure of the nervous system: neurons and their supporting cells. \textit{Oxford} (1991).

61. Bogovic, J. A., Hanslovsky, P., Wong, A. \& Saalfeld, S. Robust registration of calcium images by learned contrast synthesis. \textit{preprint; arXiv} \textbf{1511.01154 [cs.CV]} (2015).

62. Marquart, G. D., Tabor, K. M., Brown, M. \& Burgess, A. High precision registration between zebrafish brain atlases using symmetric diffeomorphic normalization. \textit{bioRxiv} \textbf{081000} (2016).


\pdfpage{21}
63. Choi, H.\textit{ et al.} Vivaldi: a domain-specific language for volume processing and visualization on distributed heterogeneous systems. \textit{IEEE Transactions on Visualization and Computer Graphics} \textbf{20}, 2407-2416, doi:10.1109/TVCG.2014.2346322 (2014).

64. Cicconet, M., Hildebrand, D. G. C. \& Elliott, H. L. To find the symmetry plane in any dimension, reflect, register, and compute a -1 eigenvector. \textit{arXiv} \textbf{1611.05971 [cs.CV]}.

65. Munkres, J. Algorithms for the assignment and transportation problems. \textit{Journal of the Society for Industrial and Applied Mathematics} \textbf{5} (1957).
\end{hangpars}


\pdfpage{22}
\section*{Acknowledgements}

We thank D.D. Bock, K.-H. Huang, and P. Huang for preliminary studies; L.-H. Ma and M.B. Ahrens for dissection help; E. Raviola, H.S. Kim, J.A. Buchanan, E.J. Benecchi, and S. Ito for tissue preparation guidance; K.J. Hayworth, J.L. Morgan, N. Kasthuri, and R. Schalek for sectioning and scanning electron microscopy advice; T. Kazimiers and J.A. Bogovic for software assistance; the Harvard Image and Data Analysis Core (H.L. Elliot, T. Xie, D.L. Richmond) and D. Berger for image processing assistance; the Harvard Center for Brain Science Neuroengineering Core (E.R. Soucy, J.S.F. Greenwood) for instrumentation support; and B.L. Shanny, A.M. Roberson, F. Gao, M.A. Afifi, F. Camacho Garcia, A.D. Wong, C.S. Elkhill, T.J. Cawley, R.J. Plummer, K.M. Runci, A. Haddad, P.E. Lewis, I. Odstrcil, A.H. Cohen, M.D. Petkova, and P.I. Petkova for reconstructions. This work was supported by the NIH through the NINDS to F.E. (DP1 NS082121, RC2 NS069407) and to W.-C.A.L. (R21 NS085320) and to the Harvard Center for Brain Science Neuroengineering Core (P30 NS062685), through the NIGMS to MMBioS via resources at the Pittsburgh Supercomputing Center (P41 GM103712), through the NCRR to the Orchestra High-Performance Compute Cluster (S10 RR028832); by the Defense Advanced Research Projects Agency (DARPA) SIMPLEX program through SPAWAR to R.B. and J.T.V. (contract N66001-15-C-4041); by the National Research Foundation of Korea (NRF) through the Bio \& Medical Technology Development Program, MSIP (NRF-2015M3A9A7029725), and the Ministry of Education Basic Science Research Program (NRF-2014R1A1A2058773) to W.-K.J.; and by training fellowships from the NIH (T32 MH20017, T32 HL007901) and NSF (IIA EAPSI award 1317014) to D.G.C.H.


\pdfpage{23}
\section*{Author Contributions}

F.E. designed the study. R.P., I.H.B., O.R., and D.G.C.H. conducted light microscopy. D.G.C.H. prepared and sectioned samples. D.G.C.H. and G.S.P. performed electron microscopy. A.W.W., S.S., and D.G.C.H. aligned electron micrographs. D.G.C.H. and O.R. completed correspondence across imaging modalities. D.G.C.H. managed neuron reconstructions. R.M.T., B.J.G., and D.G.C.H. post-processed and analysed reconstructions. M.C. and D.G.C.H. performed symmetry analyses. W.C., T.M.Q., J.M., D.G.C.H., and W.-K.J. produced visualizations. A.S.C. modified hosting and annotation software. A.B., K.L., R.B., and J.T.V. provided dataset hosting. F.E., J.W.L., W.-C.A.L., and A.F.S. managed various aspects of the project and provided resources. D.G.C.H., F.E., and J.W.L. wrote the manuscript with input from other authors.

\section*{Author Information}

The authors declare no competing financial interests. Correspondence and requests for materials should be addressed to D.G.C.H. (david@hildebrand.name) or F.E. (florian@mcb.harvard.edu).

\pdfpage{24}
\begin{center}
\includegraphics[width=0.95\linewidth,height=0.6\textheight,keepaspectratio]{figures/fig1.png}
\end{center}
\noindent\textbf{Figure 1: Targeted, multi-scale ssEM of the larval zebrafish brain.}

\noindent\textbf{a}, Volume rendering of and reslices through 16000 aligned serial micrographs that form a nearly isotropic (56.4×56.4×60nm\textsuperscript{3}vx\textsuperscript{–1}) image volume containing the anterior quarter of a larval zebrafish. \textbf{b–c}, Flattened contour reslice images reveal brain features visible with nearly isotropic resolution. \textbf{b}, The Mauthner cell (M), axon cap (AC), and axon (Ax) surrounded by other neurons, myelinated structures, and neuropil. \textbf{c}, Posterior extension of the Mauthner axon. \textbf{d}, Micrographs from a single section acquired in separate rounds of imaging. All larval zebrafish tissue in 16000 sections was acquired at nearly isotropic resolution (solid border, \textbf{e–f}). The brain alone in 12546 sections was acquired at higher (18.8×18.8×60nm\textsuperscript{3}vx\textsuperscript{–1}) resolution (dashed border, \textbf{g–h}). \textbf{e–f}, Large-caliber myelinated structures (arrowheads) are identifiable in nearly isotropic micrographs that include peripheral nerves (\textbf{e}) and sensory organs such as the ear (\textbf{f}, image corresponds to the dashed region in \textbf{d} from a nearby section). \textbf{g–h}, Many neuronal processes including myelinated fibers can be segmented from higher resolution micrographs. \textbf{i–k}, Additional rounds of targeted re-imaging can be performed to distinguish finer neuronal
\pdfpage{25}structures and the connections between them. Scale box: \textbf{a}, 50×50×50µm\textsuperscript{3}. Scale bars: \textbf{b}–\textbf{c}, 10µm; \textbf{d}, 50µm; \textbf{e}–\textbf{f}, 5µm; \textbf{g}–\textbf{h}, 1µm; \textbf{i}–\textbf{k}, 500nm.
\alttext{Composite of EM images on a white background, panels a–k. \textbf{a}: volume rendering in greys of the anterior larval zebrafish seen obliquely from the front and above: the two eyes (each a dark ring around a grey interior) sit in front, the head is cut open by a horizontal reslice (outlined in magenta) and a block-shaped reslice rises at the top left; a small yellow rectangle on the block marks the region of panel c and an orange outline marks the region of panel b. A green plane outline and a green contour line on the head lead via a green arrow down to panel d. Scale cube at lower left labelled "50 µm" with axes x, y, z. \textbf{b}: flattened contour reslice (orange frame, trapezoidal edge pointing to the orange outline in a) through brain tissue, with blue labels "Ax" (upper left, a pale band running left), "AC" (centre) and "M" (a large round pale cell body at right); surrounded by densely textured neuropil and dark myelin profiles; scale bar below (10 µm per caption). \textbf{c}: yellow-framed reslice showing a horizontal band of dark, layered (myelinated) profiles running across the middle between paler cell bodies above and paler tissue below; scale bar "10 µm". \textbf{d}: green-framed single transverse section through the head, light grey tissue on white: the brain (darker, roughly heart-shaped) in the upper middle and two dark crescent-shaped structures below it on either side; further tissue profiles below. A dashed green quadrilateral encloses the brain (the 18.8 nm px\textsuperscript{−1} region; blue rotated labels "56.4 nm px\textsuperscript{−1}" and "18.8 nm px\textsuperscript{−1}" at upper left); blue boxes mark the positions of panels h (left brain), g (centre, with a dot), e (right, outside the brain) and f (dashed box at far left). Label "z = 8185" at upper right; x/y axis arrows at lower right; scale bar "50 µm". Column headings above the right panels: "56.4 × 56.4 × 60 nm\textsuperscript{3} vx\textsuperscript{−1}" over e–f, "18.8 × 18.8 × 60 nm\textsuperscript{3} vx\textsuperscript{−1}" over g–h, "4 × 4 × 60 nm\textsuperscript{3} vx\textsuperscript{−1}" over i–k. \textbf{e}: pale nerve cross-section with an oval bundle in which four small dark ring profiles (myelinated axons) are marked by blue arrowheads; scale bar "5 µm" (for e–f). \textbf{f}: section edge with a row of cell profiles and elongated dark myelinated structures marked by three blue arrowheads at lower right; label "z = 8442" at upper right. \textbf{g}: higher-resolution neuropil with several thick dark myelin rings; a blue box labelled i encloses one of them. \textbf{h}: higher-resolution neuropil with two blue boxes labelled k (left) and j (right); scale bar "1 µm" (for g–h). \textbf{i}: a single large myelinated axon with a thick, multi-layered dark sheath around a pale axoplasm. \textbf{j}: small processes outlined by dark membranes, with a round profile enclosed in another at upper centre. \textbf{k}: fine processes outlined by dark membranes, with a small dark dot at upper centre; scale bar "500 nm" (for i–k).}

\pdfpage{26}
\begin{center}
\includegraphics[width=0.95\linewidth,height=0.6\textheight,keepaspectratio]{figures/fig2.png}
\end{center}
\noindent\textbf{Figure 2: Neuronal morphology reconstructions capturing sensory input and motor output.}

\noindent\textbf{a}, The path of a bipolar lateral line afferent neuron (orange) is tracked from a synaptic contact with a neuromast hair cell (\textbf{b}–\textbf{d}) through the peripheral lateral line ganglion (\textbf{e}) into the hindbrain, as revealed by a contour reslice through ∼5000 serial sections. Inset at top right is a lateral view of reconstructions in this figure, illustrating their position. \textbf{b}, Cross-section through the dorsal neuromast innervated by the highlighted afferent. \textbf{c}, Magnified ribbon synapse connecting a neuromast hair cell to the highlighted afferent. \textbf{d}, Series of micrographs showing the afferent’s exit from the neuromast and myelination. \textbf{e}, Cross-section through the peripheral lateral line ganglion showing the myelinated soma of the afferent. \textbf{f}, The position of a caudal primary (CaP) motor neuron in the spinal cord and its innervation of myotome 6 projected onto a reslice through ∼2200 serial sections. Scale bars: \textbf{a} top right inset, 100µm; \textbf{a}, \textbf{e}–\textbf{f}, 10µm; \textbf{b}–\textbf{d}, 1µm.
\alttext{EM composite, panels a–f, with blue labels. \textbf{a}: wide grey contour reslice (posterior at left, anterior at right; axis arrows "z x" / "y" at lower right; text "posterior ↔ anterior" below) along which an orange line labelled "lateral line afferent" runs from a neuromast at upper left, along the body surface, down and right in a U-shaped course through a dotted blue outline labelled "peripheral lateral line ganglion" at the bottom, and then up into a dotted blue outline labelled "brain" at right, where it ends. A vertical green dashed line at the neuromast connects via green lines to panel b, and one at the ganglion connects via green lines to panel e. Inset at top right: lateral view of the fish (eye at right) with the orange afferent and a red neuron, blue boxes labelled "a" and "f" marking the regions of panels a and f; small scale bar. Scale bar "10 µm" below a. \textbf{b}: transverse EM section through the neuromast (label "neuromast", "z = 4000"), a cluster of grey cells with dark irregular nuclei; the afferent profile (orange) at the bottom inside a blue box labelled c; x/y axis arrows. \textbf{c}: magnified view: a darker "hair cell" (labelled) at top with a dark round body near its lower edge, contacting the orange "afferent" profile below; a blue box labelled d at left. \textbf{d}: eight small micrographs in two rows labelled z = 4006, 4012, 4022, 4030 (top) and 4039, 4054, 4070, 4086 (bottom), each with the orange afferent profile at the centre containing darker round inclusions; the profile becomes smaller and elongated, and at 4086 it is surrounded by a dark multi-layered ring (myelin). Scale bar "1 µm" below. \textbf{e}: ganglion cross-section ("ganglion", "z = 6647", x/y axes): round cell bodies, several with dark irregular outlines, one filled orange near the centre; scale bar (10 µm per caption). \textbf{f}: reslice (posterior ↔ anterior; z/y axes) through the dorsal "spinal cord" (labelled, outlined by a blue dotted line, densely packed small profiles) and underlying muscle, with a red neuron labelled "caudal primary (CaP) motor neuron": a red dot (soma) in the spinal cord at upper centre, from which a red axon descends through a blue dotted outline labelled "myotome 6" and branches at the bottom of it. Scale bar "10 µm".}

\pdfpage{27}
\begin{center}
\includegraphics[width=0.95\linewidth,height=0.6\textheight,keepaspectratio]{figures/fig3.png}
\end{center}
\noindent\textbf{Figure 3: Reconstruction of a larval zebrafish projectome consisting of myelinated axons.}

\noindent\textbf{a–b}, Dorsal (\textbf{a}) and lateral (\textbf{b}) views of myelinated axon reconstructions manually extracted from the multi-resolution ssEM dataset. Colors are randomly assigned. \textbf{c}, Dorsal (top) and lateral (bottom) views of reconstructed lateral line primary sensory neurons. Anterior (purple) afferents were identified by their innervation of neuromasts in the periphery, where darker purple indicates more anterior. Posterior (yellow) afferents were identified by their presence in the posterior lateral line nerve. \textbf{d}, Dorsal (top) and lateral (bottom) views of a reticulospinal neuron subset including the Mauthner cell and nucleus of the medial longitudinal fasiculus (nucMLF) neurons. Cell identity was assigned by position and morphology for the Mauthner cell (purple) and the
\pdfpage{28}four canonical nucMLF neurons: MeLc (green), MeLr (yellow), MeLm (orange), and MeM (blue). \textbf{e}, Magnified view from \textbf{c} reveals topographically arranged afferent hindbrain axons such that those innervating more anterior neuromasts sit ventrolateral to those innervating more posteriorly located neuromasts. \textbf{f}, Magnified view of Mauthner cell and nucMLF neuron morphologies from \textbf{d}. Note the apparent bilateral symmetry in \textbf{c}–\textbf{f}. Scale bars: \textbf{a}–\textbf{d}, 100µm; \textbf{e}–\textbf{f}, 50µm.
\alttext{Renderings of coloured reconstructions over a semi-transparent grey volume rendering of the larval zebrafish head (anterior to the right, eyes as large dark discs with a darker round centre). \textbf{a} ("projectome consisting of myelinated axons"): dorsal view; hundreds of thin lines in many random colours run along the hindbrain/spinal region at left, cross the midline in dense bundles in the centre, and spread anteriorly between and around the two eyes; scale bar at lower left. \textbf{b}: lateral view of the same reconstructions, dense multicoloured bundles running from the left edge into the head above and behind the eye, with some lines descending below and in front of the eye; scale bar "100 µm" (for a–b). \textbf{c} ("lateral line afferents"): dorsal view (top) and lateral view (bottom); yellow afferents run in from the left (posterior) and purple/magenta afferents of various shades run from in front of and around the eyes; the two sides look mirror images of each other in the dorsal view; a blue box labelled e in the lateral view marks the region magnified in e; scale bars. \textbf{d} ("reticulospinal neurons"): dorsal (top) and lateral (bottom) views of a small set of coloured neurons: a purple Mauthner pair with dendrites lateral to the midline at the left of a blue box labelled f, and blue, green, yellow, orange and grey neurons near the front of the box; their axons run together along the midline to the left; scale bars "100 µm". \textbf{e}: magnified lateral view of the afferents in the hindbrain: yellow afferents (key at left, rotated text "posterior" in yellow, "anterior" in magenta) and magenta afferents of several shades bend through the hindbrain, the magenta ones in a thick bundle from lower right; scale bar "50 µm". \textbf{f}: magnified dorsal view: a purple Mauthner cell on each side at left (vertical branched dendrites) and, at right, green, yellow, orange, blue and grey neurons on both sides; their axons run in two parallel tracts just either side of the midline to the left edge, and the left and right sets look mirror-symmetric. Rotated colour key at left: "MeM" (blue), "MeLr" (yellow), "MeLm" (orange), "MeLc" (green), "Mauthner" (purple). Scale bar "50 µm".}
\decompnote{The caption spells "fasiculus"; elsewhere the paper writes "fasciculus".}

\pdfpage{29}
\begin{center}
\includegraphics[width=0.95\linewidth,height=0.6\textheight,keepaspectratio]{figures/fig4.png}
\end{center}
\noindent\textbf{Figure 4: Symmetry in myelinated axons of the reticulospinal system.}

\noindent\textbf{a}–\textbf{c}, Three-dimensional (3D) symmetry analysis for 44 identified neurons (22 per side) whose axon traveled in the ∼30µm-diameter medial longitudinal fasciculus (MLF). \textbf{a}, Dorsal view of identified reticulospinal axons overlaid by the co-registered Z-Brain spinal backfills label. The plane of symmetry was fit from all identified axons that formed part of the MLF and used in subsequent analyses. \textbf{b}, Pairwise bilateral symmetry comparisons in 3D. Each axon on one side was compared to the reflection of every axon on the other side across the symmetry plane. Costs were calculated as the sum of the Euclidian distances between points matched by a dynamic time warping algorithm, then normalized by the number of matches and penalized by any unmatched regions (see Extended Data Fig. 9 and Methods). The matrix is ranked by cost between left-right homologs along the diagonal (top-left to bottom-right). Asterisks indicate the results of a globally optimal pairwise assignment. Though the analysis lacked any bias based on neuron identity, in all but one pair (red asterisks), the axon on the left was assigned to its homolog on the right. Low costs observed in off-diagonal elements highlight interesting similarities between the different neuron types. \textbf{c}, Top and side views of pairs resulting in the lowest (left) and highest (right) costs. \textbf{d}–\textbf{j}, Two-dimensional (2D) symmetry analysis of neighbour relations between axon pairs in one bundle with respect to the reflected pair from the other side for 6 left-right pairs of identified neurons (12 in total). Axon identity color code in \textbf{f} applies to all \textbf{d}–\textbf{j} panels. \textbf{d}, In cross-section, most axons within the left and right MLF appear to share mirror symmetric
\pdfpage{30}relative positions. Line indicates midline. \textbf{e}, Angle and distance differences for all 15 pairs in a single cross-section (slice 5560, indicated by red line in \textbf{f}–\textbf{i}), where the vector connecting a pair on one side is compared to the reflected vector between the same pair from the opposite side (see Methods). \textbf{f}, Top view of 6 identified axon pairs for which neighbor relations were examined. \textbf{g}–\textbf{h}, Linearizing the angle (\textbf{g}) and distance (\textbf{h}) differences between all pairs reveals that positional arrangements are mirror symmetric with their contralateral counterparts over long stretches of the MLF. Additionally, the neighbour relations return to a state of symmetry even after local perturbations. Black regions indicate where reconstruction data does not exist for a given pair. \textbf{i}, Summing normalized angle and distance differences shows particular regions (peaks) with neighbour relation perturbations, such as the entry of a new axon into the bundle. Swapping the identity of one pair of axons (MeLc and Mauthner) on one side results in a consistent near-doubling of this sum over for all cross-sections in which they both exist. \textbf{j}, Calculating the angle and distance difference variance across all cross-sections identifies a few pairs with weak neighbor relations. Scale bars: \textbf{a}, 100µm; \textbf{c},\textbf{f}, 50µm; \textbf{d}, 5µm.
\alttext{\textbf{a}: dorsal volume rendering of the larval head (grey, anterior right, two dark eyes) with blue "reticulospinal reconstructions" (axons running along the midline from the left edge and branching into cell bodies and dendrites in the hindbrain and midbrain), diffuse red "Z-Brain spinal backfill" label over the same region, and a horizontal orange line along the midline labelled "plane of symmetry"; scale bar "100 µm". \textbf{b}: 22 × 22 matrix of coloured squares, rows labelled "left" and columns "right" with the same 22 neuron names in the same order (MiD3cm, RoM2L, RoM1C, MeLc, RoM3M \#1, MeLr, MiD3i, Mauthner, RoM3L, MiD2cm, MiD2i, MiR1, MiM1 \#1, MiR2, RoM2M, MiM1 \#2, MeMd, RoM3M \#2, MeLm, MeMv, CaV, RoM1R; MeLc green, MeLr yellow, Mauthner purple, MeMd light blue, MeLm orange, MeMv blue lettering). Colour key "cost" from 2000 (yellow) via orange, green and light blue to 6000 (dark navy). Most cells are dark navy (≥ 6000); the diagonal cells carry black-framed boxes with black asterisks and run from yellow at the top left (lowest costs) to green and light blue at the bottom right (CaV ≈ 3600, RoM1R ≈ 5000 by the key); in rows MiD2i and MiM1 \#1 the asterisks are red and sit off the diagonal, at MiD2i/MiM1 \#1 and MiM1 \#1/MiD2i. Off-diagonal cells of lower cost (yellow to light blue) are scattered across the matrix. \textbf{c}: two examples with top and side views of axon traces on grey backgrounds: "lowest cost: MiD3cm right vs. reflected MiD3cm left" (green and magenta traces lying on top of each other; yellow box with asterisk; "cost = 1315"; scale bar "50 µm"), and "highest cost: MeMv right vs. reflected CaV left" (green and magenta traces clearly apart in both views; dark navy square; "cost = 14608"). \textbf{d}: EM cross-section ("z = 5560") of the left and right MLF separated by a vertical blue midline; coloured dots mark the 6 axons on each side (purple, green, light blue, yellow, orange, dark blue), arranged as mirror images; scale bar "5 µm". \textbf{e} ("difference matrix at z = 5560"): 6 × 6 matrix, upper triangle labelled \textit{a}\textsubscript{z=5560}, lower triangle \textit{d}\textsubscript{z=5560}, diagonal empty; rows and columns marked with the six coloured dots; key from blue (0°/0) to green (180°/8 µm). Nearly all cells are blue (small differences); one upper cell (MeMd–MeLr) is bright green (≈ 160°) and a few are slightly lighter blue. \textbf{f}: top view of the 6 axon pairs, coloured traces converging from the right (where the neurons branch) into two tight parallel bundles running to the left; key "MeMv, MeMd, MeLr, MeLm, MeLc, Mauthner"; a red vertical line at slice 5560 extends through panels f–i (labelled "d-e" at the bottom); scale bar "50 µm". \textbf{g} ("angle difference", with a small diagram of the angle \textit{a} between two pair vectors) and \textbf{h} ("distance difference", diagram of \textit{d}): heatmaps with one row per pair (15 rows marked by pairs of coloured dots) against slice; colour keys 0°–180° and 0–8 µm; mostly blue with scattered green streaks, black blocks where data are missing (left part of several rows, a step-shaped black region at the right end, and black rows to the right for pairs that end). \textbf{i}: line plot, y "sum of all normalized \textit{a} and \textit{d}" (0 to 10), x "position in series, \textit{slice}" (labels 2013, 5013, 8013, 11013); black line "original" mostly between ≈ 0 and 4, with a sharp spike to ≈ 5 near slice 8800; pink line "swapped right [green dot] and [purple dot]" lying above the black line (≈ 2–9, maximum ≈ 9 near slice 4450) for most of the range and at or below the black line beyond slice ≈ 8900. \textbf{j} ("variance per pair across \textit{z}"): 6 × 6 matrix, upper triangle "variance in \textit{a}", lower "variance in \textit{d}"; key 0°–22.5° and 0–1 µm; mostly blue, with greener cells at MeMv–MeLm and MeMd–MeLr (upper, ≈ 18°) and a blue-green cell at MeLm–MeMv (lower, ≈ 0.55).}
\figuredata{b}{figures/fig4b.csv}{figures/fig4b.py}{digitized from the 150 ppi raster image (tier 3): cell colours matched to the printed colour key, checked against the costs printed in Fig. 4c and Extended Data Fig. 9}
\figuredata{e}{figures/fig4e.csv}{figures/fig4e.py}{digitized from the 150 ppi raster image (tier 3): cell colours matched to the printed colour key}
\figuredata{g}{figures/fig4g.csv}{figures/fig4g.py}{digitized from the 150 ppi raster image (tier 3): grid at native resolution, pixel colours matched to the printed colour key}
\figuredata{h}{figures/fig4h.csv}{figures/fig4h.py}{digitized from the 150 ppi raster image (tier 3): grid at native resolution, pixel colours matched to the printed colour key}
\figuredata{i}{figures/fig4i.csv}{figures/fig4i.py}{digitized from the 150 ppi raster image (tier 3): both lines traced per pixel column}
\figuredata{j}{figures/fig4j.csv}{figures/fig4j.py}{digitized from the 150 ppi raster image (tier 3): cell colours matched to the printed colour key}

\pdfpage{31}
\begin{center}
\includegraphics[width=0.95\linewidth,height=0.6\textheight,keepaspectratio]{figures/edfig1.png}
\end{center}
\noindent\textbf{Extended Data Figure 1: Preparing larval zebrafish brain tissue for electron microscopy.}

\noindent\textbf{a}, Immersion of intact specimen into tissue processing solutions typically resulted in poor preservation of brain ultrastructure due to a membrane diffusion barrier (arrowheads). \textbf{b–c}, Dissecting away skin and membranes permitted diffusion of solutions, resulting in improved preservation of brain ultrastructure. To minimize damage, dissections were initiated by puncturing the rhombencephalic ventricle dorsal to the hindbrain with a sharpened tungsten needle (red cross). Small anterior-directed incisions along the midline as close to the surface as possible were then made until the brain was exposed from the hindbrain to the anterior optic tectum (red dashed line). \textbf{d–f}, Following dissection and aldehyde fixation (\textbf{d}), samples were post-fixed with a reduced osmium solution (\textbf{e}) and stained with uranyl acetate (\textbf{f}). \textbf{g–h}, Processed larvae were then dehydrated with acetonitrile, infiltrated with a low-viscosity resin, mounted in a
\pdfpage{32}micromachined pre-cast resin mold to orient the sample for transverse sectioning (\textbf{g}), and surrounded by a support tissue (murine cortex) that was found to stabilize sectioning (\textbf{h}). \textbf{i}, Representative ultrastructure achieved using this process acquired as a transmission electron micrograph from a section through the optic tectum of an early dissection test specimen. Scale bars: \textbf{g}–\textbf{h}, 1mm; \textbf{d}–\textbf{f}, 500µm; \textbf{b}, 100µm; \textbf{a},\textbf{c},\textbf{i}, 1µm.
\alttext{Workflow figure. \textbf{a}: grey EM micrograph, rotated label "without dissection" at left, headings "outside brain" (left) and "inside" (right) separated by blue arrowheads at top and bottom pointing at a dark, wavy, continuous boundary running vertically just right of centre; the tissue left of it shows large pale spaces and layered membranes, the right side is fainter and less distinct; scale bar "1 µm". \textbf{b}: schematic dorsal view of a larval zebrafish head (anterior right; two black eyes, grey brain regions outlined by grey dashed lines, fins at left), title "dissection"; a red cross over the hindbrain and a red dashed arrow running anteriorly along the midline to the front of the brain; scale bar "100 µm". \textbf{c}: EM micrograph, rotated label "with dissection", headings "outside" and "inside brain" separated by blue arrowheads at a thin boundary; cells on both sides have crisp, continuous membranes; scale bar "1 µm". An arrow labelled "aldehyde fixation" leads from b to d. \textbf{d}–\textbf{f}: photographs of a whole larva from above (head right), each with scale bar "500 µm": d pale/whitish with dark eyes; arrow "reduced osmium post-fixation" to e, reddish-brown; arrow "uranyl acetate \textit{en bloc} staining" to f, dark brown to black. Arrow "dehydration, resin infiltration, and mounting" from f to g. \textbf{g}: rendering of an orange cylindrical "pre-cast resin orienting mold" with four holes, a grey "larval zebrafish" standing upright in a central hole; scale bar "1 mm". Arrow "support tissue addition and embedding" to h. \textbf{h}: photograph looking down onto an amber resin block containing a dark trapezoidal piece labelled "support tissue", in whose centre a white dashed circle encloses the "larval zebrafish"; a white reflection spot at lower left; scale bar "1 mm". \textbf{i}: wide transmission EM micrograph of neuropil: a large nucleus with granular texture and a sharply defined envelope at upper left/centre, surrounded by many profiles of neurites with crisp membranes, mitochondria with cristae, myelin-like ring profiles and bundles of fine processes; scale bar "1 µm".}

\pdfpage{33}
\begin{center}
\includegraphics[width=0.95\linewidth,height=0.6\textheight,keepaspectratio]{figures/edfig2.png}
\end{center}
\noindent\textbf{Extended Data Figure 2: Serial sectioning and ultrathin section library assembly for a 5.5dpf larval zebrafish.}

\noindent\textbf{a}, Serial sections of resin-embedded samples were picked up with an automated tape-collecting ultramicrotome modified for compatibility with larger reels containing enough tape to accommodate tens of thousands of sections. \textbf{b–c}, Direct-to-tape sectioning resulted in consistent section spacing and orientation. Just as a section left the diamond knife (blue), it was caught by the tape. \textbf{d}, After serial sectioning, the tape was divided onto silicon wafers that functioned as a stage in a scanning electron microscope and formed an ultrathin section library. For a series containing all of a 5.5dpf larval zebrafish brain, ∼68m of tape was divided onto 80 wafers (with ∼227 sections per wafer). \textbf{e}, Wafer images were used as a coarse guide for targeting electron microscopic imaging. Fiducial markers (copper circles) further provided a reference for a per-wafer coordinate system, enabling storage of the position associated with each section and, thus, multiple rounds of re-imaging at varying resolutions as needed. \textbf{f}, Low-resolution overview micrographs (758.8×758.8×60nm\textsuperscript{3}vx\textsuperscript{–1}) were acquired for each section to ascertain sectioning
\pdfpage{34}reliability and determine the extents of the ultrathin section library. Scale boxes: \textbf{a}, 5×5×5cm\textsuperscript{3}; \textbf{b}, 1×1×1cm\textsuperscript{3}; \textbf{c}, 1×1×1mm\textsuperscript{3}. Scale bars: \textbf{e}, 1cm; \textbf{f}, 250µm.
\alttext{\textbf{a}: 3D rendering (CAD style) of an ultramicrotome (cream-coloured body with binocular eyepieces at top) with a grey tape-collecting unit attached at the left: two large transparent green reels stacked vertically, rollers and guides, and an orange tape running from the reels to the knife area (blue knife boat) in front of the specimen arm. Green scale cube labelled "5 cm". An arrow downwards from a to d labelled "∼68 meters of tape divided onto silicon wafers". \textbf{b}: close-up rendering of the knife area: the orange tape, with a dashed yellow line along it, runs over a grey guide down to the blue knife boat, next to a yellow specimen block at the upper right; green scale cube "1 cm". \textbf{c}: further close-up: the orange tape emerging from the water-filled blue boat carries a row of four yellow, lozenge-shaped sections evenly spaced and in the same orientation, the newest just leaving the knife edge in front of the cylindrical yellow block face; green scale cube "1 mm". \textbf{d}: rendering of a stack of many thin grey wafers forming a cylinder, rotated label "80 wafers", with one wafer pulled out near the bottom showing the brown tape strips on its top; a purple curved arrow points from it to e. \textbf{e}: photograph of one circular silicon wafer (purple outline) covered with about 11 parallel vertical strips of dark brown tape, each carrying a column of small dark sections; handwriting on the orange labels at the top ("130201 zf M2") and bottom ("W62-130321"); three small copper circles (fiducial markers) near the edges; a small green square near the centre marks one section, connected by green lines to f; scale bar "1 cm". \textbf{f}: low-resolution overview EM image of one section (green frame): a lozenge-shaped section of support tissue with layered texture, in the centre a pale region containing the larval zebrafish cross-section (two dark oval profiles with darker round centres, and a smaller grey profile between them), enclosed by a green dashed circle; scale bar "250 µm".}

\pdfpage{35}
\begin{center}
\includegraphics[width=0.95\linewidth,height=0.6\textheight,keepaspectratio]{figures/edfig3.png}
\end{center}
\noindent\textbf{Extended Data Figure 3: Serial sectioning through the anterior quarter of a 5.5dpf larval zebrafish.}

\noindent\textbf{a}, Overview micrographs (758.8×758.8×60nm\textsuperscript{3}vx\textsuperscript{–1}) from a collection of 18207 × ∼60 nm-thick transverse serial sections that span 1.09mm through a larval zebrafish. Embedding the larval zebrafish (green dashed circle) in a support tissue (murine cortex) stabilized sectioning. Dashed lines show where overview images were cropped. \textbf{b}, Volume rendering of aligned overview micrographs. Magenta and yellow planes correspond to reslice planes in \textbf{c}. Green plane corresponds to section outlined in \textbf{a}. \textbf{c}, Reslice planes through an aligned overview image volume reveal structures contained within the series and illustrate the sectioning plane relative to the horizontal (top) and sagittal (bottom) body planes. This series spans from myotome 7 through the anterior larval zebrafish, encompassing part of the spinal cord and the entire brain. Line colours indicate the intersection of the reslice planes. \textbf{d}, Histograms of lost, partial (missing any larval zebrafish tissue), or adjacent (lost-partial or partial-partial) events per bin of 50 sections. In total, 244 (1.34\%) sections were lost and 283 (1.55\%) were partial in this series. No two adjacent sections were lost. Inset histograms expand the shaded regions to provide a detailed view of sectioning reliability with bin sizes of 5 sections. Dashed lines indicate the number of lost sections if uniformly distributed throughout the series. Scale box: \textbf{b}, 250×250×250µm\textsuperscript{3}. Scale bars: \textbf{a},\textbf{c}, 250µm.
\alttext{\textbf{a}: perspective stack of overview EM sections (grey on white), labelled z = 1, z = 2, z = 3, three dots, and z = 18207 at the back; the front section (green outline, dashed top edge) shows a lozenge of support tissue with layered texture and, inside a green dashed circle, the small cross-section of the larva in a pale region. A bracket along the stack is labelled "1.09 mm" and "posterior ↔ anterior"; x/y axes; scale bar "250 µm". \textbf{b}: volume rendering of the aligned overview stack (the larva, eye visible as a dark disc) cut by a green plane (the section in a), a yellow plane and a magenta plane; grey scale cube "250 µm". \textbf{c}: two reslices of the overview volume, posterior at left, anterior at right: top, magenta-framed "horizontal" reslice (left ↔ right; z/x axes) showing the body and head from above with both eyes (dark rims) at right, and a yellow dashed line along the length; bottom, yellow-framed "sagittal" reslice (dorsal ↔ ventral; z/y axes) showing the body from the side with an eye at the lower right, a small dark curved structure in the middle, and a magenta dashed line running obliquely; scale bar "250 µm", text "posterior ↔ anterior". \textbf{d}: three histograms on a common x axis "position in series, z" (0 to 18000, labels at 0, 6000, 12000, 18000), y label "sections per bin". "lost (244 sections, 1.34\%)", y axis to "12/50": grey bars in clusters around z ≈ 700–1700, 2600–3800, 4800–5800 (tallest bar ≈ 11 near z ≈ 5350, inside a blue shaded band), 7700–8300 and 9300–10000, plus scattered single bars from z ≈ 300 to ≈ 17900; a dashed horizontal line just above zero. Inset "2/5" with x axis 5000, 5250, 5500 (blue background) showing the 5-section bins of the shaded band, bars of height 1 or 2. "partial (283 sections, 1.55\%)", y axis to "10/50": a tall cluster at z ≈ 1300–1700 (shaded band, peak ≈ 10 near z ≈ 1500), many bars between ≈ 600 and 10000 (mostly 1–3, a few up to ≈ 6–7), and a few up to ≈ 12500; dashed line just above zero; inset (1250, 1500, 1750) with bars of 1–2. "adjacent (5 section pairs, 0.03\%)", y axis "1/50": five thin bars of height 1 near z ≈ 1000, 8250, 8450, 10000 and 12500.}
\figuredata{d}{figures/edfig3d.csv}{figures/edfig3d.py}{digitized from the 150 ppi raster image (tier 3); main histograms per pixel column (50-section bins not resolved), insets per 5-section bin}

\pdfpage{36}
\begin{center}
\includegraphics[width=0.95\linewidth,height=0.6\textheight,keepaspectratio]{figures/edfig4.png}
\end{center}
\noindent\textbf{Extended Data Figure 4: Serial sectioning through most of a 7dpf larval zebrafish brain.}

\noindent\textbf{a}, Overview micrographs from a collection of 15046 × ∼50 nm-thick transverse serial sections that span 0.75mm through a 7dpf larval zebrafish. Surrounding part of the larval zebrafish (green dashed circle) with a support tissue (murine cortex) stabilized sectioning. Dashed lines show where overview images were cropped. \textbf{b}, Volume rendering of aligned overview micrographs. Magenta and yellow planes correspond to reslice planes in \textbf{c}. Green plane corresponds to section outlined in \textbf{a}. \textbf{c}, Reslice planes through an aligned overview image volume reveal structures contained within the series and illustrate the sectioning plane relative to the horizontal (top) and sagittal (bottom) body planes. This series spans from posterior hindbrain through the anterior larval zebrafish, encompassing most of the brain. Line colours indicate the intersection of the reslice planes. \textbf{d}, Histograms depicting the number of lost, partial (missing any larval zebrafish tissue), or adjacent (lost-partial or partial-partial) events per bin of 50 sections throughout the series. In total, 6 (0.04\%) sections were lost and 25 (0.17\%) were partial in this series. No two adjacent sections were lost. Scale box: \textbf{b}, 250×250×250µm\textsuperscript{3}. Scale bars: \textbf{a},\textbf{c}, 250µm.
\alttext{\textbf{a}: perspective stack of overview EM sections labelled z = 1, z = 2, z = 3, three dots and z = 15046; the front section (green outline) shows the larva's cross-section inside a green dashed circle, sitting on top of support tissue that fills only the lower part of the section (above it is empty resin); bracket "0.75 mm", "posterior ↔ anterior"; y/x axes; scale bar "250 µm". \textbf{b}: volume rendering of the aligned stack cut by green, yellow and magenta planes; grey scale cube "250 µm". \textbf{c}: magenta-framed "horizontal" reslice (left ↔ right; z/x axes) of the head from above, brain in the middle, two eyes with dark rims at the right, a yellow dashed line running slightly obliquely; yellow-framed "sagittal" reslice (dorsal ↔ ventral; z/y axes) with an eye at the lower right and a magenta dashed oblique line; vertical banding in both; scale bar "250 µm", "posterior ↔ anterior". \textbf{d}: three histograms over "position in series, z" (0 to 15000, labels 0, 5000, 10000, 15000), y label "sections per bin". "lost (6 sections, 0.04\%)", y axis "2/50", text "6 sections": one bar of 2 near z ≈ 300 and bars of 1 near 3350, 3500, 4100 and 12300. "partial (25 sections, 0.17\%)", y axis "4/50", text "25 sections": thin bars near z ≈ 300, 1000, a cluster at ≈ 2400–3100 (tallest ≈ 4 near 2680), and bars near 3350, 3700, 4100, 8200 and 12300. "adjacent (6 section pairs, 0.04\%)", y axis "1/50", text "6 section pairs": bars of 1 near z ≈ 300, 2680, 3350, 4100 and 8200 (the bar near 300 is wider and may hold two).}
\figuredata{d}{figures/edfig4d.csv}{figures/edfig4d.py}{digitized from the 150 ppi raster image (tier 3); one row per bar run, 50-section bins resolved to about one bin}

\pdfpage{37}
\begin{center}
\includegraphics[width=0.95\linewidth,height=0.6\textheight,keepaspectratio]{figures/edfig5.png}
\end{center}
\noindent\textbf{Extended Data Figure 5: Description of partial sections.}

\noindent Collected sections were deemed partial if any larval zebrafish tissue appeared to be missing. In total, 283 sections out of 18207 were classified as partial. Those that were imaged at higher resolution were further categorized into minor, moderate, or severe subclasses. In minor cases, only tissue outside the brain was absent. Moderate cases lacked less than half of the brain. Severe cases were missing more than or equal to half of the brain. Note that it is possible that apparently missing tissue is contained in a slightly thicker adjacent section, in which case it is not entirely lost and may be accessible with different imaging strategies. \textbf{a–c}, Posterior examples of partial sections from each category. The line and arrow indicate the orientation and direction of sectioning. \textbf{e–f}, Expanded views of brain tissue from the sections depicted in \textbf{a–c}. Red dashed contours define the brain outline from an adjacent section. \textbf{g–h}, Anterior examples of partial sections from each category. \textbf{j}, Number of sections in each category for the 208 partial sections contained within the 16000 imaged at nearly isotropic (56.4×56.4×60nm\textsuperscript{3}vx\textsuperscript{–1}) resolution. Scale bars: \textbf{a}–\textbf{i}, 50µm.
\decompnote{The caption labels the expanded views "e–f" and the anterior examples "g–h"; the figure shows three expanded views (d, e, f) and three anterior examples (g, h, i), so "d–f" and "g–i" seem meant.}
\alttext{\textbf{a}–\textbf{c}: three transverse overview EM sections of the posterior larval zebrafish (light grey on white), headed "minor" (a, z = 3304), "moderate" (b, z = 3401) and "severe" (c, z = 3780). Each shows the dorsal neural tube/brain at top, a large empty oval cavity in the middle and ventral tissue with grey folded structures at the bottom, flanked by thin elongated tissue strips. A blue box at the top of each marks the region shown in d, e and f respectively. In a, a blue line with an arrow at upper left (sectioning orientation and direction). Tissue in c is visibly missing at the upper left of the box. Scale bars at lower left. \textbf{d}–\textbf{f}: enlargements of the boxed regions; a red dashed contour outlines the brain: in d the outlined region is complete (polygonal cell profiles inside), in e it is mostly present with tissue missing along the upper-left edge, and in f the upper-left part inside the contour is empty and only the lower-right part has tissue. Scale bars. \textbf{g}–\textbf{i}: three more anterior overview sections headed "minor" (g, z = 8334), "moderate" (h, z = 8332) and "severe" (i, z = 8313), each with the brain outlined by a red dashed contour at the top; in g the brain region is darker and intact, in h it is paler with some texture, in i it is very pale and largely featureless inside the contour; scale bars ("50 µm" under g). \textbf{j}: bar chart, y axis "sections" 0 to 90 (ticks 0, 30, 60, 90), three grey bars: minor ≈ 84, moderate ≈ 81, severe ≈ 43; text "208 total" at upper right.}
\figuredata{j}{figures/edfig5j.csv}{figures/edfig5j.py}{digitized from the 150 ppi raster image (tier 3)}

\pdfpage{38}
\begin{center}
\includegraphics[width=0.95\linewidth,height=0.6\textheight,keepaspectratio]{figures/edfig6.png}
\end{center}
\noindent\textbf{Extended Data Figure 6: Software modifications for seamless interaction with multiple co-registered ssEM datasets and reference atlas overlays.} Reconstructing neuronal structures across multi-resolution ssEM image volumes acquired from the same specimen profits from being able to simultaneously access and view separate but co-registered datasets. Without this ability, some of the time benefits of our imaging approach would be offset by the need to register and track structures across volumes that span both low-resolution, large fields of view and high-resolution, specific regions of interest. With this in mind, we added a feature to the Collaborative Annotation Toolkit for Massive Amounts of Image Data (CATMAID) neuronal circuit mapping software to overlay and combine image stacks acquired
\pdfpage{39}with varying resolutions in a single viewer. This feature is now available in the main CATMAID distribution. Furthermore, all datasets described here are openly accessible and hosted using this improved version of CATMAID software (http://hildebrand16.neurodata.io). \textbf{a}, Single-section images from two co-registrered ssEM datasets acquired at different resolutions from the same specimen. The combined view (left) overlays high resolution data (right, 4×4×60nm\textsuperscript{3}vx\textsuperscript{–1}) onto nearly isotropic resolution data (middle, 56.4×56.4×60nm\textsuperscript{3}vx\textsuperscript{–1}). \textbf{b}, Integrated view of multiple co-registered ssEM datasets with manual reconstructions (coloured dots) and the spinal backfill label from the Z-Brain atlas. Visible in this section is the red spinal backfill fluorescence directly on top of, for example, the Mauthner cell body (arrowhead).
\alttext{Two browser screenshots of the CATMAID web interface (address bar starting "hildebrand16.neurodata.io/catmaid/?pid=6\&zp=…", toolbars with Home, Projects, Stacks and tool icons; most interface text is too small to read). \textbf{a}: three side-by-side stack viewers showing the same EM field (cells with dark membranes at lower left, striated muscle-like tissue at right) at the same position; the right viewer's image is a slightly rotated quadrilateral on black; each viewer has a blue scale bar labelled "20 µm" (as far as legible); coordinates "[-739, 32, 8969]" at the bottom right. \textbf{b}: one viewer showing a transverse EM section of the brain (grey cell bodies and neuropil) overlaid with diffuse red/pink fluorescence patches (spinal backfill label) arranged roughly symmetrically about the midline, and many small magenta and a few yellow dots (reconstruction nodes). A blue arrowhead at right-centre points at a yellow-circled node on a red-overlaid cell body; a green tag "Mauthner\_R" at the top right of the viewer; blue scale bar "50 µm" (as far as legible) at lower left; coordinates "[7389, 1736, 8807]" at the bottom right.}
\decompnote{"co-registrered" is printed so in the caption.}

\pdfpage{40}
\begin{center}
\includegraphics[width=0.95\linewidth,height=0.6\textheight,keepaspectratio]{figures/edfig7.png}
\end{center}
\noindent\textbf{Extended Data Figure 7: Correspondence of neuron identity across whole-brain \textit{in vivo} light and \textit{post hoc} ssEM datasets.} Co-registration of \textit{in vivo} light microscopy and EM datasets can be accomplished with thin-plate spline coordinate transformations guided by manually identified landmarks. \textbf{a}, Volume renderings of the ssEM dataset (top), warped \textit{in vivo} two-photon imaging of \textit{elavl3}:GCaMP5G fluorescence from the same specimen (middle), and a merge (bottom). Reslice planes shown in \textbf{b} are indicated by magenta planes. \textbf{b}, Near-horizontal reslice planes from the ssEM volume (top) and the warped \textit{in vivo} light microscopy image volume (bottom) show gross correspondence throughout the brain across imaging modalities. \textbf{c–d}, Magnified views reveal single-neuron correspondence in the optic tectum (\textbf{c}) and telencephalon (\textbf{d}). \textbf{e–g}, This exercise revealed the imaging conditions, labelling density, and structural tissue features necessary for reliable matching across imaging modalities. It was difficult in regions (enclosed by dotted contours) where fluorescence signal was low (\textbf{e}), where many cells were packed closely together (\textbf{f}), and
\pdfpage{41}where progenitors add new neurons in between light microscopy and preparation for ssEM (\textbf{g}). Improving the light-level data with specific labelling of all nuclei and faster light-sheet or other imaging approaches should greatly improve the ease and accuracy of matching neuron identity across modalities. This ability to assign neuron identity across imaging modalities demonstrates proof-of-principle for the potential integration of rich neuronal activity maps with subsequent whole-brain structural examination of functionally characterized neurons and their networks. Arrowheads in \textbf{c} indicate the same structures as observed in each modality. Elongated structures are blood vessels. Scale box: \textbf{a}, 100×100×100µm\textsuperscript{3}. Scale bars: \textbf{b}, 100µm; \textbf{c}–\textbf{g}, 10µm.
\alttext{\textbf{a}: three volume renderings of the larval zebrafish head seen from the side (anterior right): top, the grey ssEM volume with a dark eye and a magenta reslice plane cutting obliquely through the head; middle, the green fluorescence volume alone with the same magenta plane; bottom, the merge, green brain inside the grey head. Scale cube "100 µm". \textbf{b}: magenta-framed near-horizontal reslice through the ssEM volume (top; grey brain, two black-rimmed eyes at right, anterior right) and the corresponding warped fluorescence plane (bottom, green on black, labelled "\textit{elavl3}:GCaMP5G"); the brain outlines match. Boxes labelled c (upper middle) and d (right, near the front) in both. Scale bar "100 µm". \textbf{c}: EM (top) and fluorescence (bottom) of the boxed tectal region; five arrowheads (blue on the EM, pale green on the fluorescence) mark the same small round bright features and an elongated branching structure in both. Scale bar. \textbf{d}: EM (top, pale round cell bodies in grey neuropil) and fluorescence (bottom, dark round holes in green) of the telencephalic box; scale bar "10 µm". \textbf{e}–\textbf{g}: further EM/fluorescence pairs with dashed contours (blue on EM, white on fluorescence): e, an oval region of cells where the fluorescence inside the contour is dark; f, a band of densely packed cells below a dashed line; g, a pale band of tissue enclosed between two dashed lines that is dark in the fluorescence. Scale bars ("10 µm" under e).}

\pdfpage{42}
\begin{center}
\includegraphics[width=0.95\linewidth,height=0.6\textheight,keepaspectratio]{figures/edfig8.png}
\end{center}
\noindent\textbf{Extended Data Figure 8: Registration of reference atlases to the multi-resolution ssEM dataset.} Cross-modal registration of the Z-Brain atlas and the Zebrafish Brain Browser allows for the characterization of specific domains within the ssEM dataset defined by, for example, genetically restricted labels (\textbf{a–f}) or retrograde labeling (\textbf{g–h}). \textbf{a},\textbf{c},\textbf{e},\textbf{g}, Dorsal (left) and lateral (right) views through combined, dual-volume renderings of the ssEM dataset and Z-Brain atlas image volumes of a \textit{elavl3}:H2BmRFP transgenic line (\textbf{a}), \textit{vglut2a}:GFP transgenic line (\textbf{c}), \textit{hcrt}:RFP transgenic line (\textbf{e}), and spinal backfill retrograde labeling (\textbf{g}). \textbf{b},\textbf{d},\textbf{f},\textbf{h}, Z-Brain atlas fluorescence signal from a \textit{elavl3}:H2BmRFP transgenic line (\textbf{b}), \textit{vglut2a}:GFP transgenic line (\textbf{d}), a \textit{hcrt}:RFP transgenic line (\textbf{f}), and spinal backfill retrograde labeling (\textbf{h}) overlaid onto horizontal reslice planes through the ssEM dataset. Clearly overlapping in the spinal backfill label and ssEM reslice (\textbf{h}) are the Mauthner cell and nucleus of the medial longitudinal fasciculus neurons positions, which indicate high quality of registration. Scale bars: \textbf{a}–\textbf{h}, 100µm.
\alttext{Four rows of three images (anterior to the right throughout). In each row: a dorsal volume rendering of the grey ssEM head (two dark eyes) with a coloured Z-Brain label inside it, a lateral rendering of the same, and a horizontal EM reslice with the label overlaid (x/z axis arrows at lower right); scale bars on each image ("100 µm" at the bottom left of g). \textbf{a}/\textbf{b} (red label "\textit{elavl3}:H2BmRFP"): red fills the whole brain from hindbrain to forebrain in a and the lateral view; in b red covers most of the brain on the reslice, with a darker pattern inside. \textbf{c}/\textbf{d} (green label "\textit{vglut2a}:GFP"): green fills the brain in the renderings; in d bright green patches lie in several bilaterally paired regions on the reslice. \textbf{e}/\textbf{f} (red label "\textit{hcrt}:RFP"): only small red patches: two large patches near the eyes, paired small patches and a fan-like structure in the midbrain region in the dorsal view, a few patches in the lateral view; in f only a faint small red spot near the midline on the reslice. \textbf{g}/\textbf{h} (red label "spinal backfill"): red bundles run in from the spinal cord at left and end in about a dozen paired red cell-body-sized spots in the hindbrain and midbrain; in h red spots overlie the EM reslice, symmetric about the midline.}

\pdfpage{43}
\begin{center}
\includegraphics[width=0.95\linewidth,height=0.75\textheight,keepaspectratio]{figures/edfig9.png}
\end{center}
\pdfpage{44}
\noindent\textbf{Extended Data Figure 9: Visual description and additional examples for symmetry analysis.} \textbf{a}–\textbf{c}, Three-dimensional (3D) symmetry analysis for one example pair. \textbf{a}, To compare the left MiD2cm axon to the right homolog, the left side was first reflected across the plane of symmetry (dotted line). \textbf{b}, Generating a value that represents the similarity in position and shape of the two axons was built using a dynamic time warping sequence matching approach (see Methods). \textbf{c}, Each cost was calculated as the sum of the Euclidian distances between points matched by a dynamic time warping algorithm, then normalized by the number of matches. The cost is then multiplied by a penalty factor that is proportional to the lengths of the sequences that remain unmatched (total length divided by matched length; not illustrated). \textbf{d}, One pair of myelinated axons were not assigned to their left-right homologs (Fig. 4b, red asterisks) in a globally optimal pairwise assignment for a selection of 44 identified reticulospinal neurons (22 on each side). Comparing views and costs for the assigned non-homologs (left) with unassigned left-right homologs (right) reveals similar matching costs in all cases. However, the combined cost for both non-homologs is slightly lower than that for both left-right homologs (by 174). Since the global assignment seeks to minimize the overall cost, this is likely why it assignment result grouped non-homologs over left-right homologs. The fact that the penalty accounting for the unmatched regions between the non-homologs did not increase the cost enough to prevent this assignment shows that it has a small effect in cases with short unmatched overhangs. \textbf{e}–\textbf{h}, Two-dimensional (2D) symmetry analysis of axon neighbour relations. \textbf{e}, In each 2D slice, for every pair of identified axons on one side, the same pair on opposite side was reflected across the plane of symmetry. Two metrics were then calculated to relate the original and reflected pairs: the angle difference, or dot product between the vectors connecting the pairs, and the distance difference, or the difference between the lengths of the vectors connecting the pairs. \textbf{f}, For each slice, a distance matrix was constructed from the angle and distance difference values relating for each pair. \textbf{g}–\textbf{h}, Separately linearizing (\textbf{g}) and the angle and distance differences between all pairs and then concatenating them (\textbf{h}) enables visualization of changes in these positional arrangements across slices. \textbf{i}–\textbf{k}, Extension of 2D symmetry analysis of axon neighbour relations to the same set of 44 identified reticulospinal neurons that were analyzed in 3D. \textbf{i}, Top view of the 22 identified axon pairs for which 2D analysis was performed. \textbf{j}–\textbf{k}, Viewing linearized angle (\textbf{j}) and distance (\textbf{k}) differences for tested pairs shows a generally similar trend toward stable neighbour relations as the smaller subset analysis. Most neighbour relations return to a state of symmetry even after local perturbations, such as the entry of a new axon into the bundle. Black regions indicate where reconstruction data does not exist for a given pair. Scale bars: \textbf{a},\textbf{d},\textbf{i}, 50µm; \textbf{b}, 5µm.
\decompnote{The figure fills PDF page 43; its caption is printed on page 44. The caption's "this is likely why it assignment result" is printed so.}
\alttext{\textbf{a} ("3D symmetry analysis example", "reflected MiD2cm left vs. MiD2cm right"): top and side views of two axon traces (green = reflected left, magenta = right) on grey, lying almost on top of each other; a dashed line above the top view; small boxes labelled b on both traces; yellow box with asterisk; "cost = 2054"; scale bar "50 µm". \textbf{b}: the boxed stretch enlarged (top and side), the two traces joined by short purple segments, legend "DTW mismatch" with a vector symbol; scale bar "5 µm". \textbf{c}: formula graphic: cost\textsubscript{DTW} = Σ ‖ … (the row of purple mismatch segments) … ‖ divided by \# … (the same segments), i.e. the summed lengths of the mismatch vectors over their number. \textbf{d}: "assigned non-homologs" (left column) and "unassigned left-right homologs" (right column), each with top and side views of green (reflected left) and magenta (right) traces: "reflected MiD2i left vs. MiM1 \#1 right", yellow box with red asterisk, "cost = 2242"; "reflected MiM1 \#1 left vs. MiD2i right", red asterisk, "cost = 1967"; "reflected MiD2i left vs. MiD2i right", yellow square, "cost = 2062"; "reflected MiM1 \#1 left vs. MiM1 \#1 right", yellow-orange square, "cost = 2321"; in all four the traces nearly coincide; scale bar "50 µm". \textbf{e} ("2D symmetry analysis example"): schematic: box "original" with a yellow and a purple dot joined by a black line on each side of a vertical midline; arrow to box "reflected" in which the other side's pair is mirrored onto the first; below, "comparisons": a magenta box "angle difference" showing the angle \textit{a} between the two pair vectors and a dashed magenta box "distance difference" showing the length difference \textit{d}; x/y axes. \textbf{f} ("construction of difference matrix"): a 6 × 6 coloured matrix (upper triangle \textit{a}, lower \textit{d}, diagonal white, coloured dots for the six axons, key 0°/0 to 180°/8 µm), the same pattern as Fig. 4e (mostly blue, one green cell); one upper cell framed in solid magenta and one lower cell in dashed magenta. \textbf{g} ("linearization of difference matrix"): a "pair" column of 15 dot pairs and, for the slices "slice − 1", "slice", "slice + 1", columns of 15 cells for \textit{a} and for \textit{d}, the current slice's column framed in red, with "…" between; one cell in the \textit{a} column is green. \textbf{h} ("concatenation of linearized difference matrices"): two heatmaps (\textit{a}, key 0°–180°; \textit{d}, key 0–8 µm) of the 15 pairs (dot pairs at left) over "position in series, \textit{slice}" with labels 5013 and 8013, a red vertical line near the left; mostly blue with green streaks and thin black vertical lines. \textbf{i}: top view of the 22 pairs of axons (many colours) forming two parallel bundles along the midline that spread out on the right, with a few axons running separately above and below; scale bar "50 µm". \textbf{j} ("angle difference", diagram) and \textbf{k} ("distance difference", diagram): large heatmaps with one thin row per pair over slices 2013 to ≈ 11500 (labels 2013, 5013, 8013, 11013), keys 0°–180° and 0–8 µm: mostly blue with many short green streaks; black blocks at the left of some rows and a large staircase-shaped black region on the right where rows end at different slices.}
\figuredata{f}{figures/edfig9f.csv}{figures/edfig9f.py}{digitized from the 150 ppi raster image (tier 3): cell colours matched to the printed colour key}
\figuredata{h (top, a)}{figures/edfig9h_a.csv}{figures/edfig9h.py}{digitized from the 150 ppi raster image (tier 3): grid at native resolution, colours matched to the printed key}
\figuredata{h (bottom, d)}{figures/edfig9h_d.csv}{figures/edfig9h.py}{digitized from the 150 ppi raster image (tier 3): grid at native resolution, colours matched to the printed key}
\figuredata{j}{figures/edfig9j.csv}{figures/edfig9j.py}{digitized from the 150 ppi raster image (tier 3): grid at native resolution; single pair rows not resolved}
\figuredata{k}{figures/edfig9k.csv}{figures/edfig9k.py}{digitized from the 150 ppi raster image (tier 3): grid at native resolution; single pair rows not resolved}

\pdfpage{45}
\begin{center}
\includegraphics[width=0.95\linewidth,height=0.6\textheight,keepaspectratio]{figures/edfig10.png}
\end{center}
\noindent\textbf{Extended Data Figure 10: Examples of non-neuronal tissues contained within the dataset.}

\noindent In addition to capturing the whole brain, the nearly isotropic (56.4×56.4×60nm\textsuperscript{3}vx\textsuperscript{–1}) image volume containing the anterior quarter of a larval zebrafish serves as a high-resolution reference atlas for several other tissues and structures. Three selected sections (\textbf{a},\textbf{h},\textbf{m}) are accompanied by example images (\textbf{b–g},\textbf{i–l},\textbf{n–o}) to illustrate the variety of tissues and structures contained within the dataset. Scale bars: \textbf{a},\textbf{h},\textbf{m}, 50µm; \textbf{i}–\textbf{l},\textbf{n}–\textbf{o}, 10µm; \textbf{b}–\textbf{g}, 5µm.
\alttext{Three transverse EM sections (left column) with blue boxes marking enlarged regions shown to their right; each enlargement carries a rotated tissue label. \textbf{a} ("z = 6133"): section through the posterior head/anterior trunk, with the dark neural tissue at the top, a central cavity below it and gut and body-wall tissues further down; boxes b–g. \textbf{b} "somitic muscle": striated muscle fibres with regular rows of small dots. \textbf{c} "pronephric duct": cells around a lumen with brush-border texture. \textbf{d} "esophagus": a fan of fine pale lines between epithelial cells. \textbf{e} "vasculature": several elongated grey oval cells with dark nuclei on a pale background; scale bar "5 µm". \textbf{f} "pectoral fin": rows of round pale cells within a stack of tissue layers. \textbf{g} "liver": cells with round nuclei and dark granular deposits. Scale bars in b–g (5 µm per caption). \textbf{h} ("z = 9257"): section through the head with the brain (dark) at the top, two small dark oval bodies at left and right, and boxes i–l below the brain. \textbf{i} "utricle": a dark, ring-shaped, layered oval body above a sensory epithelium with elongated cells. \textbf{j} "chondrocranium": cartilage with scattered cells in pale matrix and an elongated oval structure at right. \textbf{k} "pharynx": epithelium with many small dark rod-like profiles at its surface; scale bar "10 µm". \textbf{l} "heart": muscle cells with dark outlines and striated texture. \textbf{m} ("z = 14066"): section through the anterior head with both eyes (large black lenses, retina with a dark pigmented rim) on either side of the forebrain; boxes n (top of the head) and o (between the eye and the brain). \textbf{n} "skin": thin surface layers over a row of pale cells; scale bar "10 µm". \textbf{o} "extraocular muscle": black pigmented retinal layer with comb-like extensions at left, muscle and pale tissue to the right. Scale bars "50 µm" under a, h and m.}

\pdfpage{46}
\textbf{Supplementary Video 1:} Survey of aligned section overview electron micrographs. Several sections excluded for faster viewing. Scale bar: 250µm.

\textbf{Supplementary Video 2:} Volume rendering of aligned section overview electron micrographs. Magenta and yellow planes correspond to reslice planes in Extended Data Figure 3c. Green plane corresponds to section outlined in Extended Data Figure 3a. Scale box: 250µm.

\textbf{Supplementary Video 3:} Survey of traverse larval zebrafish electron micrographs throughout the section series after image co-registration. Several sections excluded for faster viewing. Scale bar: 50µm.

\textbf{Supplementary Video 4:} Flythrough of serial micrographs centered on a reconstruction of the left Mauthner cell (magenta dot) beginning from its lateral dendrite, through the soma, through the axon cap, crossing to the contralateral side, and ending at its position in the spinal cord. Vertical line segments through these micrographs correspond to the contour reslices in Figure 1a−c. Several sections excluded for faster viewing. Scale bar: 5µm.

\textbf{Supplementary Video 5:} Matched horizontal reslice planes through the ssEM dataset and warped \textit{in vivo} two-photon e\textit{lavl3}:GCaMP5G fluorescence datasets. Scale bar: 100µm.

\textbf{Supplementary Video 6:} Serial electron micrographs centered on a reconstruction of a lateral line afferent neuron innervating the right dorsal (D2) neuromast (orange dot; Figure 3a−e). The video starts in high-resolution (4×4×60nm\textsuperscript{3}vx\textsuperscript{–1}) data of the neuromast where the afferent receives input from a hair cell (Figure 3b−c). The field-of-view then expands to follow the myelinated afferent in nearly isotropic resolution (56.4×56.4×60nm\textsuperscript{3}vx\textsuperscript{–1}) data throughout the periphery. It then magnifies to higher resolution (18.8×18.8×60nm\textsuperscript{3}vx\textsuperscript{–1}) data when approaching the afferent soma position in the peripheral lateral line ganglion before the afferent unmyelinates and enters the brain. Scale bar: 1µm (size changes depending on zoom in video).
\decompnote{The lateral line afferent, its neuromast and hair-cell input are shown in Figure 2a–e and 2b–c; Figure 3 is the projectome. "Figure 3a−e" and "Figure 3b−c" here seem to mean Figure 2.}

\textbf{Supplementary Video 7:} Rotating volume rendering of the projectome consisting of myelinated axons manually reconstructed from the multi-resolution ssEM dataset. Colors are randomly assigned. Scale bar: 100µm.

\textbf{Supplementary Video 8:} Rotating rendering of lateral line system reconstructions revealed a striking bilateral symmetry. Posterior lateral line afferents are shown in yellow. Anterior and cranial lateral line afferents are shown in purple. Most neuromasts (blue) have a bilaterally symmetric counterpart except in one case (orange). Scale bar: 100µm.

\textbf{Supplementary Video 9:} Volume rendering of ssEM data and the Z-Brain spinal backfill label (red) combined with reticulospinal neuron reconstructions. This cross-modal registration of Z-Brain atlas data allowed for the anatomical characterization of domains within the ssEM dataset and aided with validation of neuron identification. Scale bar: 100µm.

\textbf{Supplementary Video 10:} Visualization of two-dimensional (2D) symmetry analysis of axon neighbour relations. Angle and distance differences (top right) were calculated from cross-sections (top left) of a subset of myelinated axon reconstructions (bottom), linearized, and
\pdfpage{47}aggregated (middle) for display. This video explains how the 2D symmetry analysis was performed across slices and illustrates how panels in Figure 4 and Extended Data Figure 9 were constructed.

\end{document}
